\documentclass[11pt,a4paper]{article}
\usepackage[utf8]{inputenc}
\usepackage[T1]{fontenc}
\usepackage[ngerman,english]{babel}
\usepackage{geometry}
\geometry{a4paper, left=2.5cm, right=2.5cm, top=2.5cm, bottom=2.5cm}
\usepackage{mathptmx}
\usepackage{helvet}
\usepackage{microtype}
\usepackage{amsmath}
\usepackage{titlesec}
\usepackage{booktabs}
\usepackage{tabularx}
\usepackage{xcolor}
\usepackage{authblk}
\usepackage{xurl}
\usepackage{enumitem}
\usepackage{graphicx}
\usepackage{tikz}
\usetikzlibrary{shapes.geometric, arrows.meta, positioning, calc, fit}
\usepackage{float}
\usepackage{setspace}
\usepackage{natbib}
\usepackage{longtable}
\usepackage{quoting}
\usepackage{multirow}
\usepackage{array}
\usepackage{hyperref}
\newcolumntype{L}[1]{>{\raggedright\arraybackslash}p{#1}}
\newcolumntype{Y}{>{\raggedright\arraybackslash}X}
\newcolumntype{C}{>{\centering\arraybackslash}X}
\newenvironment{CompactLongTable}{%
  \begingroup
  \footnotesize
  \setstretch{1.0}%
  \setlength{\LTpre}{0pt}%
  \setlength{\LTpost}{0pt}%
}{%
  \endgroup
}

\titleformat{\section}{\Large\bfseries\sffamily\color{black}}{\thesection}{1em}{}
\titleformat{\subsection}{\large\bfseries\sffamily\color{darkgray}}{\thesubsection}{1em}{}
\titleformat{\subsubsection}{\normalsize\bfseries\sffamily\color{darkgray}}{\thesubsubsection}{1em}{}

\hypersetup{
    pdftitle={Ein integriertes multiaxiales Modell zur computergestützten psychiatrischen Diagnostik},
    pdfauthor={Lukas Geiger},
    pdfsubject={Methoden- und Designpapier für ein 6-Achsen-Entscheidungsunterstützungssystem},
    pdfkeywords={Multiaxiale Diagnostik, DSM-5-TR, ICD-11, ICF, Entscheidungsunterstützung, Abdeckungsanalyse},
    unicode=true,
    colorlinks=true,
    linkcolor=black,
    urlcolor=blue,
    citecolor=black,
    breaklinks=true,
    hypertexnames=false
}
\makeatletter
\renewcommand*{\HyperDestNameFilter}[1]{ger.#1}
\makeatother
\setlength{\emergencystretch}{4em}

\onehalfspacing

\begin{document}

% ===================================================================
% TITELSEITE
% ===================================================================
\selectlanguage{ngerman}

\title{\textbf{\huge Ein integriertes multiaxiales Modell\\zur computergest\"utzten psychiatrischen Diagnostik}\\[0.5em]
\Large Synthese von DSM-5-TR, ICD-11 und ICF\\in einem 6-Achsen-Entscheidungsunterst\"utzungssystem\\[0.3em]
\large Ein Methoden- und Designpapier}

\author[1]{Lukas Geiger\thanks{Korrespondenz: Lukas Geiger, Bernau, Deutschland. ORCID: \href{https://orcid.org/0009-0005-7296-1534}{0009-0005-7296-1534}}}
\affil[1]{Unabh\"angiger Forscher, Bernau}

\date{Juni 2026}

\maketitle

\begin{abstract}
\noindent Die Abschaffung des multiaxialen Diagnosesystems im DSM-5 (2013) hinterlie\ss{} eine strukturelle L\"ucke in der psychiatrischen Diagnostik. Die vorliegende Arbeit pr\"asentiert ein 6-Achsen-Modell zur computergest\"utzten multiaxialen Diagnostik, das die kategoriale Klassifikation nach DSM-5-TR und ICD-11 mit der funktionalen Erfassung der ICF in einem Entscheidungsunterst\"utzungssystem vereint. Das Modell umfasst sechs Achsen: Achse~I (Psychische Profile mit temporaler Diagnostik und quantitativer Abdeckungsanalyse), Achse~II (dimensionale Pers\"onlichkeitsdiagnostik mittels PID-5 und biographische Kontextualisierung), Achse~III (Medizinische Synopse in symmetrischer Parallelstruktur zu Achse~I), Achse~IV (ICF-Integration, WHODAS~2.0, Cultural Formulation), Achse~V (operationalisierte Fallformulierung im 3P/4P-Modell) und Achse~VI (Evidenz-Matrix, CAVE-Warnhinweise, longitudinale Symptomverlaufsdokumentation). Drei zentrale Designbeitr\"age werden hervorgehoben: (1)~die quantitative Abdeckungsanalyse (Diagnostic Coverage Index) als prozentbasierte Metrik f\"ur diagnostische Vollst\"andigkeit, (2)~die symmetrische Achse-I/III-Architektur f\"ur interprofessionelle Zusammenarbeit und (3)~die operationalisierte Verbindung von Klassifikation und individualisierter Fallformulierung. Die 6-Stufen-Gatekeeper-Logik folgt der von First (2024) beschriebenen Differenzialdiagnostik-Sequenz. Eine Python-Referenzimplementierung als hierarchische Zustandsmaschine ist \"offentlich verf\"ugbar. Dies ist ein Methoden- und Designpapier; die empirische Validierung steht noch aus.

\vspace{0.5em}
\noindent \textbf{Schl\"usselbegriffe:} Multiaxiale Diagnostik, DSM-5-TR, ICD-11, Abdeckungsanalyse, Entscheidungsunterst\"utzungssystem, HiTOP, HL7 FHIR, Fallformulierung, ICF
\end{abstract}

\newpage
\begingroup
\footnotesize
\setstretch{1.0}
\setlength{\parskip}{0pt}
\tableofcontents
\endgroup

% ===================================================================
% KI-NUTZUNGSERKLAERUNG (NEUTRAL, NICHT-AUTORSCHAFTLICH)
% ===================================================================
\section*{Erkl\"arung zur Nutzung KI-gest\"utzter Werkzeuge}
\addcontentsline{toc}{section}{Erkl\"arung zur Nutzung KI-gest\"utzter Werkzeuge}

Generative KI-Werkzeuge wurden als nicht-autorschaftliche Hilfsmittel f\"ur Literaturorganisation, Textstrukturierung, \"Ubersetzungsunterst\"utzung, sprachliche \"Uberarbeitung, Konsistenzpr\"ufungen und LaTeX/PDF-Qualit\"atssicherung eingesetzt. Ihre Ausgaben wurden als vorl\"aufige Vorschl\"age behandelt und vom Autor gepr\"uft, \"uberarbeitet oder verworfen. Sie dienten weder als Autoren oder Koautoren noch als Danksagungsempf\"anger, Fakteninstanzen oder unabh\"angige Validierungsinstanzen. Forschungsdesign, Quellenauswahl, wissenschaftliche Argumentation, klinische Interpretation, finale redaktionelle Entscheidungen und Verantwortung f\"ur das Manuskript liegen ausschlie\ss{}lich beim Autor.

\newpage

% ===================================================================
% METHODIK
% ===================================================================
\section{Methodik}

Diese Arbeit folgt einer Design-Science-Methodik f\"ur die Entwicklung eines klinischen Artefakts zur Entscheidungsunterst\"utzung. Der Forschungsprozess umfasste vier Phasen:

\textbf{Phase~1: Problemidentifikation und Anforderungsanalyse.} Eine strukturierte Literatur\"ubersicht zur Abschaffung des DSM-IV-Multiaxialsystems identifizierte f\"unf spezifische klinische Defizite: Verlust der strukturierten Aufforderung zur biopsychosozialen Beurteilung, geringe GAF-Interrater-Reliabilit\"at, inkonsistente Achse-IV-Nutzung, die k\"unstliche Achse-I/II-Grenze und die niedrige Adoptionsrate der Ersatz-V/Z-Codes. Anforderungen wurden aus diesen Defiziten, aus Firsts (2024) Goldstandard-Rahmen f\"ur Differenzialdiagnostik und aus aktuellen dimensionalen Ans\"atzen (HiTOP, RDoC, ICD-11 dimensionales Pers\"onlichkeitsmodell) abgeleitet.

\textbf{Phase~2: Iteratives Design.} Die 6-Achsen-Architektur wurde durch iterative Zyklen von Design und Verfeinerung entwickelt. Designentscheidungen --- wie die symmetrische Achse-I/III-Struktur, das Konzept der Abdeckungsanalyse und das CAVE-Warnsystem --- wurden gegen die identifizierten Anforderungen evaluiert. Jede Iteration wurde dokumentiert und versioniert (V1--V5). Es ist anzumerken, dass der Designprozess kein formales Expertenpanel einschloss; die iterative Verfeinerung erfolgte prim\"ar durch den Autor mit KI-gest\"utzter Konsistenzpr\"ufung. Eine formale Expertenbegutachtung ist als Phase~1 der Validierungsstrategie (Abschnitt~\ref{sec:validierung}) vorgesehen.

\textbf{Phase~3: Prototypische Implementierung.} Eine Python-Referenzimplementierung (ca.\ 1.850~Zeilen) unter Verwendung hierarchischer Zustandsmaschinen (\texttt{transitions}-Bibliothek) und einer Streamlit-Oberfl\"ache wurde entwickelt, um die Machbarkeit zu demonstrieren. Der Quellcode ist im \href{https://github.com/um-bruch/multiaxial-diagnostic-system}{Projekt-Repository auf GitHub} \"offentlich verf\"ugbar.

\textbf{Phase~4: KI-gest\"utzte Prozessunterst\"utzung.} Der Autor nutzte generative KI-Werkzeuge als nicht-autorschaftliche Hilfsmittel f\"ur Literaturorganisation, Textstrukturierung, sprachliche \"Uberarbeitung und Konsistenzpr\"ufung. Alle generierten Vorschl\"age unterlagen der Pr\"ufung durch den Autor; die KI-gest\"utzte Arbeit ersetzte weder empirische Validierung noch Expertenbegutachtung. Die wissenschaftliche Verantwortung f\"ur alle Inhalte liegt beim Autor.

\textbf{Umfang und Limitationen der vorliegenden Arbeit.} Dies ist ein Methoden- und Designpapier. Es pr\"asentiert Architektur, Begr\"undung und prototypische Implementierung des 6-Achsen-Modells. Die empirische Validierung (Interrater-Reliabilit\"at, klinischer Nutzen, konvergente Validit\"at) ist in vier Phasen geplant (Abschnitt~\ref{sec:validierung}), aber noch nicht durchgef\"uhrt. Die Abdeckungsanalyse-Schwellenwerte (85\%/60\%) sind als initiale Werte auf Basis klinischer Erw\"agungen vorgeschlagen und bed\"urfen empirischer Kalibrierung in zuk\"unftigen Studien. Gem\"a\ss{} der Design-Science-Methodik \citep{Hevner2004, Peffers2007} steht die formale Evaluationsphase des Artefakts noch aus; die in Abschnitt~\ref{sec:validierung} beschriebene Validierungsstrategie beschreibt den geplanten Evaluationsrahmen.

\begin{table}[ht]
\centering
\caption{Vergleich computergest\"utzter diagnostischer Systeme: OPCRIT \citep{McGuffin1991}, CATEGO \citep{Wing1974}, DTREE \citep{First1993}, CIDI \citep{Kessler2004}, MHIRA \citep{Stasiak2023}. DTREE- und CIDI-Eintr\"age basieren auf ver\"offentlichter Dokumentation. Zu beachten: Alle bestehenden Systeme verf\"ugen \"uber empirische Validierung, die f\"ur das vorliegende Modell noch aussteht.}
\scriptsize
\setlength{\tabcolsep}{2pt}
\begin{tabularx}{\textwidth}{@{}Y *{6}{C}@{}}
\toprule
Merkmal & OPCRIT & CATEGO & DTREE & CIDI & MHIRA & 6-Achsen-\\Modell \\
\midrule
Multiaxiale Struktur & Nein & Nein & Nein & Nein & Nein & Ja (6 Achsen) \\
Duale Kodierung (DSM/ICD) & Partiell & Nur ICD & DSM/ICD & DSM/ICD & Flexibel & Voll\-st\"andig \\
Abdeckungs-\\analyse & Nein & Nein & Nein & Nein & Nein & Ja \\
Dimensionale Integration & Nein & Nein & Nein & Nein & Partiell & Ja\\(HiTOP) \\
FHIR-Interoperabi\-lit\"at & Nein & Nein & Nein & Nein & Partiell & Geplant \\
Empirische Validierung & Ja & Ja & Ja & Ja (ext.) & Partiell & Geplant \\
\bottomrule
\end{tabularx}
\setlength{\tabcolsep}{6pt}
\normalsize
\label{tab:vergleich}
\end{table}

\newpage

% ===================================================================
% TEIL I: GRUNDLAGEN UND PROBLEMSTELLUNG
% ===================================================================

\section{Einleitung: Die diagnostische L\"ucke nach Abschaffung des multiaxialen Systems}
\label{sec:einleitung}

Die Abschaffung des multiaxialen Diagnosesystems im DSM-5 \citep{APA2013} hinterlie\ss{} eine strukturelle L\"ucke in der psychiatrischen Diagnostik. Das urspr\"ungliche, mit dem DSM-III \citep{APA1980DSM3} eingef\"uhrte System umfasste f\"unf Achsen: Achse~I (klinische St\"orungen), Achse~II (Pers\"onlichkeitsst\"orungen und geistige Behinderung), Achse~III (medizinische Krankheitsfaktoren), Achse~IV (psychosoziale und umweltbezogene Probleme in 9~Kategorien) und Achse~V (Globale Beurteilung des Funktionsniveaus, GAF-Score 0--100). Die APA eliminierte dieses System nach einem bereits 2004 initiierten Abschaffungsantrag aus folgenden Gr\"unden: geringe Interrater-Reliabilit\"at des GAF-Scores, konzeptuelle Vermengung von Symptomen und Funktionsniveau in Achse~V, inkonsistente klinische Nutzung von Achse~IV und die k\"unstliche Grenze zwischen Achse~I und~II.

Die klinische Gemeinschaft verlor dabei mehr als sie gewann. \citet{Probst2014} dokumentierte, dass die Eliminierung von Achse~IV \glqq die Notwendigkeit der Kontextber\"ucksichtigung nicht eliminiert\grqq{}. \citet{Kress2014} fanden, dass Kliniker nun \glqq umso wachsamer systematische Wege zur Erfassung biopsychosozialer Informationen\grqq{} finden m\"ussen --- ohne die strukturierten Aufforderungen des alten Systems. Die als Achse-IV-Ersatz eingef\"uhrten V/Z-Codes verzeichnen eine geringe Adoptionsrate: \citet{ErlichFirst2025} konstatieren, dass \glqq ein allgemeines Bewusstsein f\"ur diese Codes und die Bedeutung ihrer Nutzung zur Kommunikation sozialer Determinanten der Gesundheit nicht eingetreten ist\grqq{}.

Am kritischsten ist, dass die Abschaffung zum Verlust der strukturierten multiaxialen Diagnostik f\"uhrte und Kliniker ohne ein systematisches Rahmenwerk zur Integration biologischer, psychologischer und sozialer Dimensionen zur\"ucklie\ss{}.

\textbf{Bestehende multiaxiale und integrative Ans\"atze.} Das vorliegende Modell baut auf mehreren grundlegenden Ans\"atzen zur multiaxialen und integrativen psychiatrischen Diagnostik auf, geht aber substanziell \"uber diese hinaus. \citet{Williams1985} dokumentierte die Urspr\"unge und fr\"uhen Kritiken des DSM-III-Multiaxialsystems und identifizierte sowohl dessen klinischen Nutzen als strukturierte Aufforderung zur umfassenden Beurteilung als auch seine operationalen Schw\"achen --- insbesondere den unzuverl\"assigen GAF-Score und die inkonsistente Anwendung der psychosozialen Stressoren-Achse. Das vorliegende Modell adressiert diese spezifischen Kritiken: Es ersetzt den GAF durch validierte Instrumente (WHODAS~2.0, ICF Core Sets), eliminiert die k\"unstliche Achse-I/II-Grenze und f\"uhrt die quantitative Abdeckungsanalyse als objektive Erg\"anzung zum klinischen Urteil ein.

\citet{Bolton2008} argumentierte, dass Fallformulierung --- der individualisierte narrative Bericht \"uber die Schwierigkeiten eines Patienten --- die kategoriale Diagnose erg\"anzen, nicht ersetzen sollte. Boltons zentrale Einsicht war, dass Klassifikationssysteme erfassen, \textit{was} ein Patient hat, w\"ahrend die Formulierung erfasst, \textit{warum} und \textit{wie}. Das vorliegende Modell operationalisiert diese Einsicht durch Achse~V (Integriertes Bedingungsmodell), die ein strukturiertes 3P/4P-Formulierungsschema mit expliziten R\"uckbez\"ugen auf die diagnostischen Achsen implementiert --- und damit die L\"ucke zwischen Boltons narrativer Formulierung und der strukturierten Klassifikation \"uberbr\"uckt, die klinische Systeme ben\"otigen.

\citet{Wakefield1992} schlug die einflussreiche \glqq Harmful Dysfunction\grqq{}-Analyse vor, wonach psychische St\"orungen Zust\"ande sind, in denen ein interner Mechanismus seine nat\"urliche Funktion nicht erf\"ullt \textit{und} die Dysfunktion dem Individuum Schaden zuf\"ugt. Dieser konzeptuelle Rahmen pr\"agt den Ansatz des vorliegenden Modells zur Differenzialdiagnostik: Die 6-Stufen-Gatekeeper-Logik (nach First, 2024) unterscheidet systematisch normale Variation von Dysfunktion, indem sie Evidenz sowohl f\"ur beeintr\"achtigte Mechanismen \textit{als auch} f\"ur klinisch bedeutsame Belastung oder Funktionsbeeintr\"achtigung verlangt. Die Abdeckungsanalyse operationalisiert dies weiter, indem sie quantifiziert, welche Symptome durch etablierte Diagnosen erkl\"art werden und welche unerkl\"art bleiben --- eine direkte Operationalisierung von Wakefields Forderung, genuine Dysfunktion statt blo\ss{}er statistischer Abweichung zu identifizieren.

\textbf{Abgrenzung zum OPD-System.} Das Operationalisierte Psychodynamische Diagnostik-System (OPD-3) stellt das wichtigste multiaxiale Diagnosesystem im deutschsprachigen Raum dar und verwendet ebenfalls einen Mehr-Achsen-Ansatz mit f\"unf Achsen (Krankheitserleben und Behandlungsvoraussetzungen, Beziehung, Konflikt, Struktur, psychische und psychosomatische St\"orungen). W\"ahrend das OPD prim\"ar psychodynamisch fundiert ist und die Therapieplanung innerhalb dieses Paradigmas unterst\"utzt, verfolgt das vorliegende Modell einen schulen\"ubergreifenden, kategorial-dimensionalen Ansatz mit Fokus auf computergest\"utzte Entscheidungsunterst\"utzung. Die Systeme sind komplement\"ar: Das OPD erfasst intrapsychische Dynamik (Konflikte, Struktur, Beziehungsmuster) in einer Tiefe, die das vorliegende System nicht anstrebt; das 6-Achsen-Modell bietet daf\"ur die quantitative Abdeckungsanalyse, die somatisch-psychische Parallelstruktur und die algorithmische Gatekeeper-Logik, die im OPD nicht enthalten sind. In der klinischen Praxis k\"onnen OPD-Befunde (insbesondere Konflikt- und Strukturachse) als Material f\"ur Achse~II (Grundkonflikte) und Achse~V (Fallformulierung) des vorliegenden Modells verwendet werden.

\textbf{Verh\"altnis zum AMDP-System.} Das AMDP-System (Arbeitsgemeinschaft f\"ur Methodik und Dokumentation in der Psychiatrie) ist im deutschsprachigen Raum das Standardinstrument f\"ur die psychopathologische Befunderhebung. Im vorliegenden Modell liefert der AMDP-Befund das \textit{Rohmaterial} f\"ur Achse~I: Die dort dokumentierten Symptome werden in der Abdeckungsanalyse (Ii) den diagnostischen Kriterien zugeordnet. Das 6-Achsen-System ersetzt den AMDP-Befund nicht, sondern baut auf ihm auf.

Bezeichnenderweise bleibt die multiaxiale Klassifikation in der Kinder- und Jugendpsychiatrie explizit erhalten: \citet{BJPsychAdv2024} argumentieren, dass das multiaxiale Klassifikationssystem ein \glqq klinisch n\"utzliches biopsychosoziales Rahmenmodell\grqq{} darstellt, das in der Erwachsenenpsychiatrie voreilig aufgegeben wurde. Diese Persistenz in einem Bereich, der inh\"arent multidisziplin\"are Zusammenarbeit erfordert, unterst\"utzt die Wiedereinf\"uhrung multiaxialer Strukturen auch f\"ur Erwachsene.

Ein berechtigter Einwand gegen die Wiedereinf\"uhrung multiaxialer Systeme lautet, dass die Probleme des DSM-IV nicht nur der Implementierung, sondern dem multiaxialen Ansatz selbst inh\"arent sein k\"onnten: Die erzwungene Aufteilung in diskrete Achsen fragmentiert m\"oglicherweise klinische Information, die besser holistisch erfasst w\"urde. Das vorliegende Modell adressiert diesen Einwand durch das Skelett-Prinzip (keine Pflichtfelder), die achsen\"ubergreifende Fallformulierung (Achse~V) und die bewusste Flexibilit\"at der interprofessionellen Zuordnung. Ob diese Designentscheidungen das Fragmentierungsproblem tats\"achlich l\"osen, ist eine empirische Frage, die erst durch die geplante Validierung beantwortet werden kann.

Die vorliegende Arbeit stellt ein 6-Achsen-Modell vor, das die spezifischen Schw\"achen des historischen Systems adressiert und zugleich F\"ahigkeiten vorschl\"agt, die \"uber bestehende Systeme hinausgehen.

\subsection{Institutionelle Konvergenz}

Das vorgeschlagene Modell steht nicht isoliert, sondern antizipiert die Richtung, in die sich die psychiatrischen Institutionen selbst bewegen. Das 17-k\"opfige \textit{APA Future DSM Strategic Committee} entwickelt derzeit einen \glqq strategischen Fahrplan\grqq{} mit vier Unterkomitees, die Dimensionalit\"at, Biomarker, Funktion/Lebensqualit\"at und sozio\"okonomische/kulturelle/umweltbezogene Determinanten explorieren. \citet{ErlichFirst2025} fordern explizit die R\"uckkehr zu einem biaxialen System zur Erfassung sozialer Determinanten der Gesundheit.


\subsection{Designphilosophie: Skelett, lebendes Dokument, maximale Freiheit}
\label{sec:designphilosophie}

Drei Designprinzipien leiten das gesamte System:

\begin{description}[style=nextline, leftmargin=1.5cm, font=\bfseries]
\item[Skelett-Prinzip] Das System fungiert als diagnostisches \textit{Skelett}: Es bietet f\"ur jede klinisch relevante Information einen definierten Ort, ohne dass irgendein Feld als Pflichtfeld deklariert wird. Kein Diagnostiker wird gezwungen, Abschnitte auszuf\"ullen, die f\"ur den aktuellen Fall irrelevant sind. Das System ist ein Korsett zur St\"utze --- nicht zur Einschr\"ankung --- des diagnostischen Prozesses.
\item[Lebendes Dokument] Die diagnostische Akte ist kein statisches Formular, sondern ein \textit{lebendes Dokument}, das mit dem klinischen Prozess w\"achst. Neue Testergebnisse, Laborberichte, Beobachtungsprotokolle und Bewertungen k\"onnen jederzeit integriert werden. Jede Achse ist offen f\"ur iterative Erg\"anzung: Ein Laborbericht wird in die Evidenz-Matrix (Achse~VI) eingetragen mit Datum, Bezeichnung, klinischer Bewertung und Achsenverweis; eine sp\"atere Fremdanamnese erg\"anzt das Kontaktprotokoll; ein neuer psychometrischer Test aktualisiert Achse~I.
\item[Bericht als Momentaufnahme] Berichte --- ob Entlassbriefe, Gutachten oder Verlaufsberichte --- sind \textit{eingefrorene Momentaufnahmen} des lebenden Dokuments zu einem bestimmten Zeitpunkt. Das Dokument selbst lebt weiter und kann nach einem Bericht neue Informationen aufnehmen, w\"ahrend der Bericht den Stand zum Zeitpunkt seiner Erstellung repr\"asentiert.
\end{description}

Diese Philosophie unterscheidet das System fundamental von starren Formularsystemen: Es bietet Struktur, wo Struktur n\"utzlich ist, und Freiheit, wo Freiheit klinisch geboten ist.

\textbf{Beitrag des vorliegenden Artikels:} Der Hauptbeitrag liegt in drei vorgeschlagenen Designelementen: (1)~der quantitativen Abdeckungsanalyse als klinische Qualit\"atssicherungsmetrik f\"ur diagnostische Vollst\"andigkeit, (2)~der symmetrischen Achse-I/III-Architektur, die strukturierte interprofessionelle Zusammenarbeit erm\"oglicht, und (3)~der operationalisierten Verbindung von kategorialer Diagnostik und individualisierter Fallformulierung in einem einzigen System. Verwandte Ans\"atze existieren in Teilbereichen (z.\,B.\ Bayesianische Diagnosesysteme, OPD-Multiaxialit\"at); der spezifische Beitrag liegt in deren Integration in ein computergest\"utztes Gesamtsystem. Diese Designelemente adressieren klinische Defizite, die seit der Abschaffung des DSM-IV-Multiaxialsystems bestehen und durch aktuelle Forderungen \citep{ErlichFirst2025} erneut Dringlichkeit erhalten haben.


% ===================================================================
% 2. DAS 6-ACHSEN-MODELL
% ===================================================================
\section{Architektur des 6-Achsen-Modells}
\label{sec:modell}

Das vorgeschlagene Modell integriert die kategoriale Diagnostik (DSM-5-TR \citep{APA2022}; ICD-11 \citep{WHO2019}) in ein epistemologisches Gesamtsystem. Es dient der Erfassung der diagnostischen Lebensspanne, der Kausalit\"atspr\"ufung zwischen Somatik und Psyche sowie der Dokumentation der Behandlungs- und Remissionsgeschichte. Ein zentrales Designprinzip ist die \textbf{symmetrische Parallelstruktur} zwischen Achse~I und Achse~III: Beide Achsen verf\"ugen \"uber analoge Subachsen (Verdachtsdiagnosen, Widerlegungen, Behandlungsgeschichte, Abdeckungsanalyse, Untersuchungsplan), sodass Psychologen und Mediziner mit identischen strukturellen Werkzeugen in ihren jeweiligen Dom\"anen arbeiten k\"onnen. Tabelle~\ref{tab:achsen} gibt eine \"Ubersicht.

\begin{figure}[ht]
\centering
\begin{tikzpicture}[
    axisbox/.style={draw, rounded corners=4pt, minimum width=3.8cm, minimum height=0.7cm,
        font=\sffamily\small\bfseries, text=white, align=center},
    subbox/.style={draw, rounded corners=2pt, minimum width=3.6cm, font=\scriptsize,
        align=left, text width=3.4cm, inner sep=3pt, fill=white},
    >=Stealth
]
% --- Achse I (links) ---
\node[axisbox, fill=blue!70!black] (ax1) at (0,0) {Achse I\\[-1pt]Psychisch};
\node[subbox, below=0.3cm of ax1, draw=blue!70!black] (ax1sub) {
    Ia: Akute St\"orungen\\
    Ib: Chronische Verl\"aufe\\
    Ic: Widerlegte Verdachte\\
    Id: Remittierte Diagnosen\\
    Ie: Remissionsfaktoren\\
    If: Behandlungsgeschichte\\
    Ig: Therapietreue\\
    Ih: Verdachtsdiagnosen\\
    Ii: Abdeckungsanalyse\\
    Ij: Untersuchungsplan
};

% --- Achse II (Mitte-links) ---
\node[axisbox, fill=teal!80!black] (ax2) at (4.5,0) {Achse II\\[-1pt]Biographie \& Pers.};
\node[subbox, below=0.3cm of ax2, draw=teal!80!black] (ax2sub) {
    PID-5-Trait-Profil\\
    Pr\"agende Erfahrungen\\
    Grundkonflikte
};

% --- Achse III (rechts von I, symmetrisch) ---
\node[axisbox, fill=red!70!black] (ax3) at (9,0) {Achse III\\[-1pt]Somatisch};
\node[subbox, below=0.3cm of ax3, draw=red!70!black] (ax3sub) {
    IIIa: Akute Diagnosen\\
    IIIb: Chron. Diagnosen (erkl.)\\
    IIIc: Beitragende Faktoren\\
    IIId: Remittierte Erkrank.\\
    IIIe: Remissionsfaktoren\\
    IIIf: Behandlungsgeschichte\\
    IIIg: Medikamentenadh\"arenz\\
    IIIh: Verdachtsdiagnosen\\
    IIIi: Abdeckungsanalyse\\
    IIIj: Untersuchungsplan\\
    IIIk: Kausalanalyse\\
    IIIl: Genetik/Familie\\
    IIIm: Medikamentenanamnese
};

% --- Symmetrie-Pfeil zwischen I und III ---
\draw[dashed, thick, blue!50!red, {Stealth[length=5pt]}-{Stealth[length=5pt]}]
    (ax1sub.east) -- (ax3sub.west)
    node[midway, above, font=\scriptsize\itshape, text=black] {symmetrische Parallelstruktur};

% --- Achse IV (unten links) ---
\node[axisbox, fill=orange!80!black] (ax4) at (0,-5.8) {Achse IV\\[-1pt]Funktionsstatus};
\node[subbox, below=0.3cm of ax4, draw=orange!80!black] (ax4sub) {
    WHODAS 2.0\\
    GdB-Einstufung\\
    ICF Core Sets\\
    Cultural Formulation (CFI)\\
    Behandlungsnetzwerk
};

% --- Achse V (unten Mitte) ---
\node[axisbox, fill=purple!70!black] (ax5) at (4.5,-5.8) {Achse V\\[-1pt]Fallformulierung};
\node[subbox, below=0.3cm of ax5, draw=purple!70!black] (ax5sub) {
    Pr\"adisponierende Faktoren\\
    Ausl\"osende Faktoren\\
    Aufrechterhaltende Faktoren\\
    Protektive Faktoren (4P)
};

% --- Achse VI (unten rechts) ---
\node[axisbox, fill=gray!70!black] (ax6) at (9,-5.8) {Achse VI\\[-1pt]Evidenz \& Sicherheit};
\node[subbox, below=0.3cm of ax6, draw=gray!70!black] (ax6sub) {
    Evidenz-Matrix\\
    CAVE-Warnungen\\
    Symptomverlauf\\
    Kontaktprotokoll
};

\end{tikzpicture}
\caption{Architektur des 6-Achsen-Modells mit Sub-Achsen. Achse~I und~III teilen eine symmetrische Parallelstruktur (angezeigt durch den gestrichelten Pfeil).}
\label{fig:architecture}
\end{figure}

\begin{table}[H]
\centering
\caption{\"Ubersicht der diagnostischen Achsen des 6-Achsen-Modells}
\label{tab:achsen}
\scriptsize
\begin{tabularx}{\textwidth}{@{}L{0.8cm}L{2.1cm}L{3.3cm}Y L{2.2cm}@{}}
\toprule
\textbf{Achse} & \textbf{Bezeichnung} & \textbf{Unterteilung} & \textbf{Kerninhalte} & \textbf{Typische Profession} \\
\midrule
I & Psychische Profile & Ia: Akute St\"orungen \newline Ib: Chronische Verl\"aufe \newline Ic: Widerlegte Verdachte \newline Id: Remittierte Diagnosen \newline Ie: Remissionsfaktoren \newline If: Behandlungsgeschichte \newline Ig: Therapietreue \newline Ih: Verdachtsdiagnosen \newline Ii: Abdeckungsanalyse \newline Ij: Untersuchungsplan & DSM-5 Cross-Cutting; PHQ-9, PCL-5, ITQ; Differenzialdiagnostik; Chronifizierung; Zeitliche Verortung & Psychologe / Psychiater \\
\addlinespace
II & Biographie \& Pers\"onlichkeit & Entwicklungshistorie \newline Pr\"agende Erfahrungen \newline Grundkonflikte \newline Pers\"onlichkeit (PID-5) & IQ, Bildung, Sozialisation; Lebensereignisse; OPD-Konflikte; PID-5-BF/+M & Psychologe / Sozialp\"adagoge \\
\addlinespace
III & Medizinische Synopse & IIIa: Akute med. Diagnosen \newline IIIb: Chron. Diagnosen (erkl.) \newline IIIc: Beitragende Faktoren \newline IIId: Remittierte Erkrank. \newline IIIe: Remissionsfaktoren \newline IIIf: Behandlungsgesch. \newline IIIg: Medikamentenadh. \newline IIIh: Verdachtsdiagnosen \newline IIIi: Abdeckungsanalyse \newline IIIj: Untersuchungsplan \newline IIIk: Kausalit\"atsanalyse \newline IIIl: Genetik/Familie \newline IIIm: Medikamenten\-anamnese & Biographische Medizin\-geschichte; Kausal\-analyse Soma--Psyche; erbbiologische Belastung; Pharmako\-therapie mit Wirkungs\-bewertung & Mediziner / Facharzt \\
\addlinespace
IV & Umwelt \& Funktion & Teilhabe (ICF) \newline Psychosoziale Umst\"ande \newline Bezugspersonen-Netzwerk \newline Cultural Formulation & WHODAS~2.0, GdB; Mini-ICF-APP; Stressoren; Behandlungsnetzwerk; CFI & Sozialarbeiter / Sozialp\"adagoge \\
\addlinespace
V & Bedingungsmodell & Pr\"adisponierende Fakt. \newline Ausl\"osende Faktoren \newline Aufrechterhaltende Fakt. \newline Protektive Faktoren \newline Klinische Synthese & 3P/4P-Modell; Achsen\-integration; Behandlungs\-planung; Prognose & Interdisziplin\"ares\\Team \\
\addlinespace
VI & Belegsammlung \& Klinische Sicherheit & Evidenz-Matrix (m. Bewertung) \newline CAVE-Warnhinweise \newline Symptomverlauf \newline Kontaktprotokoll & Dokumentenregister m. klin. Bewertung; CAVE-Warnhinweise; Symptomverlauf; Kontakt-/Beobachtungs\-protokoll & Alle Professionen \\
\bottomrule
\end{tabularx}
\normalsize
\end{table}

\subsection{Achse I: Psychische Profile --- Temporale Diagnostik}
\label{sec:achse1}

Achse~I geht weit \"uber die einfache Diagnoseliste des DSM-5 hinaus, indem sie Diagnosen in ihrer zeitlichen Dynamik erfasst. Die Unterteilung in zehn Subachsen (Ia--Ij) erm\"oglicht die Dokumentation des gesamten diagnostischen Lebenslaufs: aktuelle akute St\"orungen (Ia), chronische Verl\"aufe mit Chronifizierungsrisiko (Ib), explizit widerlegte Verdachtsdiagnosen mit differenzialdiagnostischer Begr\"undung (Ic), remittierte Diagnosen mit belegter Inaktivit\"at (Id), Remissionsfaktoren wie Therapie, Medikation oder spontane Bew\"altigung (Ie), vollst\"andige Behandlungsgeschichte einschlie\ss{}lich Wirkungen und Nebenwirkungen (If), Therapietreue aus Selbst- und Fremdperspektive (Ig), laufende Verdachtsdiagnosen als diagnostische Hypothesen (Ih), die Abdeckungsanalyse unerkl\"arter Restsymptomatik (Ii) und strukturierte Untersuchungspl\"ane f\"ur n\"achste diagnostische Schritte (Ij).

Die Reihenfolge von Abdeckungsanalyse (Ii) vor Untersuchungsplan (Ij) folgt der klinischen Logik: Erst m\"ussen diagnostisch unerkl\"arte Symptome systematisch identifiziert werden, bevor gezielte Untersuchungen zu deren Kl\"arung geplant werden k\"onnen. Der Untersuchungsplan ergibt sich somit direkt aus den L\"ucken der Abdeckungsanalyse.

\textbf{PRO/CONTRA-Evidenzbewertung:} Jede Diagnose erh\"alt eine strukturierte Evidenzbewertung mit expliziter Dokumentation der Befunde, die f\"ur (PRO) und gegen (CONTRA) die Diagnose sprechen, sowie eine numerische Konfidenzsch\"atzung (0--100\%). Diese Operationalisierung folgt dem Prinzip der differenzialdiagnostischen Transparenz: Der Kliniker muss nicht nur dokumentieren, \textit{dass} eine Diagnose gestellt wird, sondern \textit{warum} --- und welche Gegenargumente er erw\"agt und verworfen hat.

\textbf{Quantitative Abdeckungsanalyse (Ii):} Eine Symptom-Diagnosen-Matrix quantifiziert den Grad der diagnostischen Erkl\"arung jedes Symptoms durch die aktiven Diagnosen. Die formale Definition, Schwellenwerte und Sensitivit\"atsanalyse werden in Abschnitt~\ref{sec:abdeckung} ausf\"uhrlich dargestellt.

\textbf{Priorisierter Untersuchungsplan (Ij):} Der Untersuchungsplan verwendet ein 3-Stufen-Priorisierungsschema: \textit{Dringend} (innerhalb von 4~Wochen abzukl\"aren --- z.\,B.\ Suizidalit\"atsabkl\"arung, medizinische Notfalldiagnostik), \textit{Wichtig} (innerhalb von 3~Monaten --- z.\,B.\ neuropsychologische Testung, Differenzialdiagnostik) und \textit{Verlaufskontrolle} (kontinuierliches Monitoring --- z.\,B.\ Symptomverlauf unter Therapie). Jede geplante Untersuchung wird einem Fachgebiet zugeordnet und mit einer klinischen Begr\"undung versehen.

Diese temporale Dimensionalit\"at fehlte sowohl in Achse~I des DSM-IV als auch im flachen Listing des DSM-5.

\subsection{Achse II: Biographie und dimensionale Pers\"onlichkeit}

Achse~II modernisiert die alte Pers\"onlichkeitsachse durch Integration der dimensionalen Trait-Diagnostik mittels PID-5, wie es sowohl ICD-11 als auch das Alternative Modell des DSM-5 empfehlen. Das ICD-11-Modell ersetzt die kategorialen Pers\"onlichkeitsst\"orungs-Typen durch ein Schweregrad-Rating (Pers\"onlichkeitsschwierigkeit $\to$ leicht $\to$ moderat $\to$ schwer), f\"unf Trait-Qualifikatoren (Negative Affektivit\"at, Distanziertheit, Dissozialit\"at, Enthemmung, Anankastie) und einen Borderline-Muster-Spezifikator. Das DSM-5-AMPD verwendet die Level of Personality Functioning Scale (LPFS) f\"ur Kriterium~A und f\"unf Trait-Dom\"anen mit 25 Facetten f\"ur Kriterium~B.

Der \textbf{PID-5-BF+M} (36~Items, 6~Dom\"anen, 18~Facetten) verbr\"uckt beide Systeme und ist in \"uber 12~Sprachen verf\"ugbar. Metaanalysen zeigen eine Gesamt-Konvergenz von $r = 0{,}62$ zwischen den Systemen, mit Dom\"anen-Korrelationen von $r = 0{,}78$--$0{,}86$, mit Ausnahme von Anankastie ($r = 0{,}34$), die kein direktes AMPD-\"Aquivalent besitzt.

Neuere Arbeiten zur Integration der dimensionalen Systeme von ICD-11 und DSM-5 f\"ur Pers\"onlichkeitsst\"orungen in eine einheitliche Taxonomie mit nicht-\"uberlappenden Traits unterst\"utzen den hier verfolgten PID-5-basierten Ansatz \citep{Bach2022}.

\"Uber die dimensionale Pers\"onlichkeitsdiagnostik (IIa) hinaus integriert Achse~II zwei weitere biographische Komponenten: \textbf{Pr\"agende Erfahrungen (IIb)} dokumentieren lebensgeschichtlich bedeutsame Ereignisse mit Lebensabschnitt und Auswirkung auf die aktuelle Situation (z.\,B.\ fr\"uhe Verlusterfahrung, Gewalterfahrung, Migration). \textbf{Grundkonflikte (IIc)} erfassen wiederkehrende intrapsychische Themen (z.\,B.\ Autonomie-Abh\"angigkeits-Konflikt, Selbstwertkonflikt), wie sie in der psychodynamischen Tradition und im OPD-System beschrieben werden. Beide Komponenten sind als offene Listen implementiert, die mit dem diagnostischen Prozess wachsen, und liefern wesentliches Material f\"ur die Fallformulierung in Achse~V.

\subsection{Achse III: Medizinische Synopse --- Symmetrische Parallelstruktur zu Achse~I}
\label{sec:achse3}

Ein zentrales Designprinzip des Modells ist die \textbf{strukturelle Symmetrie} zwischen Achse~I (psychisch) und Achse~III (somatisch). Die Einsicht dahinter: Wenn ein Psychologe in Achse~I Verdachtsdiagnosen formulieren, Widerlegungen dokumentieren und Abdeckungsanalysen durchf\"uhren kann, muss ein Mediziner in Achse~III \"uber exakt dieselben M\"oglichkeiten verf\"ugen. Nur so funktioniert das multi-professionelle Modell.

Achse~III umfasst daher 13~Subachsen:

\begin{description}[style=nextline, leftmargin=1.5cm, font=\bfseries]
\item[IIIa--IIIj: Parallele Kernstruktur] Diese zehn Subachsen folgen der Grundstruktur von Achse~I, wobei die Subachsen IIIb (erkl\"arende chronische Diagnosen) und IIIc (beitragende Faktoren) eine dom\"anenspezifische Differenzierung darstellen, die in Achse~I kein direktes \"Aquivalent hat: akute somatische Diagnosen (IIIa), chronische somatische Diagnosen mit vollst\"andig erkl\"arendem Charakter (IIIb), beitragende medizinische Faktoren --- somatische Befunde, die Psychopathologie beg\"unstigen oder verst\"arken, ohne sie vollst\"andig zu erkl\"aren (IIIc), remittierte somatische Erkrankungen (IIId), somatische Remissionsfaktoren --- operative Eingriffe, Medikation, Lebensstil\"anderung (IIIe), medizinische Behandlungsgeschichte einschlie\ss{}lich Operationen, Pharmakotherapie und Nebenwirkungen (IIIf), Medikamentenadh\"arenz aus \"arztlicher Dokumentation und Patientenbericht (IIIg), medizinische Verdachtsdiagnosen als laufende diagnostische Hypothesen (IIIh), somatische Abdeckungsanalyse --- K\"orpersymptome, die durch keine aktuelle medizinische Diagnose erkl\"art werden (IIIi), und medizinischer Untersuchungsplan f\"ur weiterf\"uhrende somatische Diagnostik (IIIj).
\item[IIIk--IIIm: Achse-III-spezifische Subachsen] Drei Subachsen haben keine Entsprechung in Achse~I, da sie die spezifische Soma--Psyche-Beziehung abbilden: integrierte Kausalit\"atsanalyse --- systematische Zuordnung somatischer Befunde zur Psychopathologie unter Differenzierung von vollst\"andiger Erkl\"arung (z.\,B.\ Hypothyreose erkl\"art Depression) und beg\"unstigender Wirkung (z.\,B.\ chronischer Schmerz verst\"arkt Angstst\"orung) (IIIk); genetische und famili\"are Belastung --- Familienanamnese f\"ur somatische und psychische Erkrankungen, bekannte genetische Pr\"adispositionen und ggf.\ Ergebnisse aus Exom-Sequenzierung (IIIl); Medikamentenanamnese und Interaktionen --- strukturierte Erfassung aller aktuellen und vergangenen Medikamente mit Dosierung, Einheit, Zweck/Indikation, Einnahmeschema, Wirkungsbewertung (0--10), Nebenwirkungen und potentiellen Interaktionen (IIIm).
\end{description}

\textbf{Anmerkung zur Strukturentscheidung IIIb/IIIc:} Die explizite Differenzierung zwischen vollst\"andig erkl\"arenden chronischen Diagnosen (IIIb) und blo\ss{} beitragenden Faktoren (IIIc) integriert die Kausalanalyse direkt in die Klassifikationsstruktur. Widerlegte medizinische Verdachtsdiagnosen werden innerhalb von IIIh (Verdachtsdiagnosen) durch einen Statuswechsel dokumentiert, analog zur klinischen Praxis, in der Verdachtsdiagnosen ausgeschlossen und mit Begr\"undung verworfen werden.

\textbf{Anmerkung zu IIIm:} Die Medikamentenanamnese als eigenst\"andige Subachse (statt blo\ss{}er Auflistung in IIIf) erm\"oglicht eine strukturierte Erfassung von Wirkungsprofilen, Interaktionsrisiken und Therapietreue auf Pr\"aparatebene. Dies ist klinisch essentiell, da Polypharmazie und Medikamenteninteraktionen eine h\"aufige Ursache sowohl f\"ur somatische als auch f\"ur psychiatrische Komplikationen darstellen.

Diese Symmetrie stellt sicher, dass das System als \textit{gemeinsames Werkzeug} f\"ur Mediziner und Psychologen funktioniert. Ein Mediziner, der Achse~III ausf\"ullt, verf\"ugt \"uber dieselbe diagnostische Tiefe wie ein Psychologe in Achse~I.

\subsection{Achse IV: Umwelt und Funktion --- ICF-Integration und Cultural Formulation}
\label{sec:achse4}

Achse~IV kombiniert mehrere validierte Instrumente. Der \textbf{WHODAS~2.0} \citep{WHO2010} (WHO Disability Assessment Schedule) erfasst sechs Dom\"anen --- Kognition, Mobilit\"at, Selbstversorgung, Umgang mit Menschen, Lebensaktivit\"aten und Partizipation --- wahlweise als 12-Item-Kurzform (erkl\"art 81\% der Varianz; \citealp{Luciano2010}) oder 36-Item-Vollversion. Die psychometrischen Eigenschaften sind stark: Cronbachs $\alpha = 0{,}94$--$0{,}96$, Test-Retest-ICC $= 0{,}93$--$0{,}96$.

WHODAS~2.0 und GAF messen \textit{fundamental unterschiedliche Konstrukte}: GAF vermengt Symptome und Funktionsniveau; WHODAS~2.0 misst Behinderung unabh\"angig von der Symptomschwere. \citet{Gspandl2018} fanden \textit{keine signifikante Korrelation} zwischen selbstbewertetem WHODAS~2.0 und GAF bei Schizophrenie-Spektrum-St\"orungen. Das System verwendet prim\"ar den WHODAS~2.0 f\"ur die Behinderungserfassung. Der GAF kann optional f\"ur die R\"uckw\"artskompatibilit\"at beibehalten werden, ist jedoch keine Kernkomponente des Modells.

F\"ur den deutschen Kontext ist die \textbf{GdB-Integration} (Grad der Behinderung) essentiell. Der GdB verwendet eine 20--100-Skala ($\geq 50$ = Schwerbehinderung) und wird nach den Versorgungsmedizinischen Grunds\"atzen (VMG) bewertet.

Drei \textbf{ICF-Core-Set}-Familien f\"ur psychische Gesundheit wurden f\"ur Depression, bipolare St\"orung und Schizophrenie entwickelt \citep{Cieza2004Depression, AyusoMateos2013Bipolar, GomezBenito2018Schizophrenia}: Depression (31~kurze/121~umfassende Kategorien), bipolare St\"orung (19/38) und Schizophrenie (25/97). F\"ur die praktische Implementierung ist das \textbf{Mini-ICF-P}-Rating von \citet{LindenBaron2005} mit 13~Kapazit\"atsdimensionen das effizienteste Werkzeug.

\textbf{Operationalisierung der ICF-Integration in Achsen~III und~VI:} Die ICF-Codierung im 6-Achsen-Modell folgt einer dreistufigen Logik. \textit{Erstens} werden auf Achse~III (Medizinische Synopse) die K\"orperfunktionen (ICF-Kapitel~b) und K\"orperstrukturen (ICF-Kapitel~s) dokumentiert, die f\"ur die somatisch-psychische Kausalanalyse (IIIk) relevant sind --- beispielsweise \texttt{b130} (Funktionen der psychischen Energie und des Antriebs), \texttt{b152} (emotionale Funktionen) oder \texttt{b144} (Funktionen des Ged\"achtnisses). \textit{Zweitens} erfasst Achse~IV die Aktivit\"aten und Partizipation (ICF-Kapitel~d) sowie Umweltfaktoren (ICF-Kapitel~e) --- operationalisiert \"uber WHODAS~2.0 (Dom\"anen d1--d6) und Mini-ICF-APP (13~Kapazit\"atsdimensionen, z.\,B.\ F\"ahigkeit zur Regelm\"a\ss{}igkeit, Flexibilit\"at, Kontaktf\"ahigkeit). \textit{Drittens} werden in der Evidenz-Matrix (Achse~VI) die ICF-Qualifikatoren (Ausma\ss{} der Beeintr\"achtigung: 0--4) zu jeder codierten Funktionseinschr\"ankung dokumentiert, um die Quantifizierung \"uber die Zeit zu erm\"oglichen.

\textbf{Mehrwert gegen\"uber dem WHODAS~2.0 allein:} W\"ahrend der WHODAS~2.0 als globales Behinderungsinstrument sechs Dom\"anen auf einer Summenskala erfasst, erm\"oglicht das 6-Achsen-Modell eine \textit{achsenbezogene Differenzierung}: Funktionseinschr\"ankungen werden nicht nur quantifiziert, sondern kausal den Achsen~I (psychisch) und~III (somatisch) zugeordnet. Die Abdeckungsanalyse (Ii/IIIi) identifiziert zudem Funktionsdefizite, die durch keine aktuelle Diagnose erkl\"art werden --- eine Dimension, die WHODAS~2.0 nicht adressiert. Das Mini-ICF-APP erg\"anzt den WHODAS~2.0 durch st\"orungsspezifische Kapazit\"atsprofile, die f\"ur die Behandlungsplanung und Prognose in Achse~V (Bedingungsmodell) direkt verwertbar sind.

\subsubsection{Bezugspersonen und Behandlungsnetzwerk}

Achse~IV erfasst neben den psychosozialen Stressoren und Funktionsinstrumenten auch das \textbf{Bezugspersonen- und Behandlungsnetzwerk} des Patienten. F\"ur jede Kontaktperson werden Name, Rolle (Elternteil, Partner, Hausarzt, Facharzt, Therapeut, Sozialarbeiter, Betreuer etc.), Institution, Kontaktdaten und Anmerkungen dokumentiert. Diese strukturierte Netzwerkerfassung dient mehreren Zwecken: Sie macht das soziale St\"utzsystem sichtbar, das f\"ur die Fallformulierung (Achse~V, protektive Faktoren) und die Behandlungsplanung essentiell ist; sie dokumentiert das beteiligte professionelle Netzwerk f\"ur die Koordination; und sie identifiziert potentielle Informationsquellen f\"ur Fremdanamnesen (dokumentiert im Kontaktprotokoll, Achse~VI).

\subsubsection{Cultural Formulation Interview (CFI)}

Das DSM-5-TR enth\"alt ein strukturiertes \textbf{Cultural Formulation Interview} (CFI) mit 16~Kernfragen in vier Dom\"anen: kulturelle Definition des Problems, kulturelle Wahrnehmung von Ursache/Kontext/Unterst\"utzung, kulturelle Faktoren in der Bew\"altigung und kulturelle Faktoren in der Kliniker-Patient-Beziehung. F\"ur ein Modell, das umfassende biopsychosoziale Erfassung beansprucht, ist die Integration des CFI als optionales Modul in Achse~IV essentiell. Die 12~supplement\"aren Module des CFI (f\"ur spezifische Populationen und Kontexte) k\"onnen bedarfsorientiert aktiviert werden.

\subsection{Achse V: Integriertes Bedingungsmodell --- Operationalisierte Fallformulierung}
\label{sec:achse5}

Achse~V stellt eine vorgeschlagene Erweiterung dar, die in bisherigen Systemen nicht enthalten ist. Das pr\"adisponierende--ausl\"osende--aufrechterhaltende Faktorenmodell (klinisches 3P/4P-Modell) verbr\"uckt diagnostische Klassifikation mit Fallformulierung --- etwas, das keine DSM-Edition jemals enthalten hat. \citet{Owen2023} argumentiert, dass dieser Formulierungsansatz \glqq die Auswahl der Fakten diszipliniert und die Behandlung zielgerichtet macht\grqq{}. Die Synthese aller Achsen --- wie Biologie (III), Biographie (II) und aktuelle Stressoren (IV) mit der Psychopathologie (I) interagieren --- wird hier expliziert.

Das Modell operationalisiert die Fallformulierung durch vier strukturierte Komponenten:

\begin{description}[style=nextline, leftmargin=1.5cm, font=\bfseries]
\item[Pr\"adisponierende Faktoren (P1)] Vulnerabilit\"atsfaktoren, die vor der St\"orung existierten. Quellen: Achse~II (biographische Risikofaktoren, Pers\"onlichkeits-Traits), Achse~III (genetische Belastung, IIIm), Achse~IV (fr\"uhe psychosoziale Belastungen). Kodierung: Faktor, Quellenachse, Evidenzst\"arke (gesichert/wahrscheinlich/m\"oglich), Beleg (Achse~VI-Referenz).
\item[Ausl\"osende Faktoren (P2)] Ereignisse oder Ver\"anderungen, die das Auftreten der St\"orung ausgel\"ost haben. Quellen: Achse~IV (aktuelle Stressoren, Life Events), Achse~III (somatische Ausl\"oser). Kodierung: Faktor, zeitlicher Zusammenhang, Quellenachse, Beleg.
\item[Aufrechterhaltende Faktoren (P3)] Bedingungen, die die St\"orung perpetuieren. Quellen: Achse~I (Komorbidit\"aten, Therapietreue), Achse~III (unbehandelte somatische Befunde), Achse~IV (fortbestehende psychosoziale Belastungen). Kodierung: Faktor, Mechanismus, \"Anderbarkeit (modifizierbar/stabil), Priorit\"at f\"ur Intervention.
\item[Protektive Faktoren (P4)] Ressourcen, die Resilienz f\"ordern. Quellen: Achse~II (Pers\"onlichkeitsst\"arken), Achse~IV (soziale Unterst\"utzung, \"okonomische Ressourcen), Achse~I (Remissionsfaktoren, Ie). Kodierung: Faktor, Quellenachse, Aktivierbarkeit.
\end{description}

\textbf{Schulen\"ubergreifende Kompatibilit\"at:} Das 3P/4P-Modell ist bewusst schulen\"ubergreifend konzipiert. Verhaltenstherapeutische Fallkonzeptionen (z.\,B.\ das SORKC-Modell) lassen sich direkt in die Achse-V-Struktur \"ubersetzen: Stimulus-Situationen entsprechen ausl\"osenden Faktoren (P2), Organismus-Variablen den pr\"adisponierenden Faktoren (P1), Reaktionsmuster und Konsequenzen den aufrechterhaltenden Faktoren (P3). Psychodynamische Formulierungen (OPD-Konflikte, strukturelle Einschr\"ankungen) finden ihren Platz ebenfalls in P1 und P3. Diese Kompatibilit\"at ist kein Zufall, sondern reflektiert die empirische Konvergenz verschiedener Therapieschulen auf gemeinsame \"atiologische Faktoren.

Die Verkn\"upfung mit der Abdeckungsanalyse (Achse~Ii) ist zentral: Symptome, die durch keine Achse-I-Diagnose abgedeckt sind, \textit{sollten} durch die Achse-V-Formulierung erkl\"arbar sein. Symptome, die weder diagnostisch noch formulatorisch erkl\"art werden, repr\"asentieren genuine diagnostische L\"ucken und l\"osen den Untersuchungsplan (Achse~Ij) aus.

\subsection{Achse VI: Belegsammlung, CAVE-Warnhinweise und Symptomverlauf}

Achse~VI wurde gegen\"uber dem Originalentwurf substanziell erweitert und umfasst nun drei Komponenten:

\begin{description}[style=nextline, leftmargin=1.5cm, font=\bfseries]
\item[Evidenz-Matrix] Ein zentrales Dokumentenverzeichnis, das jede kodierte Information mit einem Beleg verkn\"upft (Dokumentenanalyse, Befragung, Testung). Jeder Eintrag umfasst neben Achsenbezug, Dokumententyp und Beschreibung ein explizites \textbf{Bewertungsfeld} f\"ur die klinische Einsch\"atzung: Ein eingehender Laborbericht wird mit Datum, Bezeichnung, Quelle \textit{und} klinischer Bewertung dokumentiert --- so ist transparent, wie der Kliniker den Befund im diagnostischen Kontext interpretiert. Dieses Designmerkmal ist in bestehenden diagnostischen Klassifikationssystemen nach Kenntnis des Autors nicht \"ublich, wobei eine systematische Literaturrecherche hierzu noch aussteht. Es unterst\"utzt direkt die klinische Rechenschaftspflicht.
\item[CAVE-Warnhinweise] Ein strukturiertes Warnsystem f\"ur kritische klinische Risiken, die achsen\"ubergreifend Relevanz haben. Jeder Warnhinweis wird mit einer Kategorie versehen: \textit{Medikamenten-Interaktion} (z.\,B.\ Lithium + NSAR), \textit{Labor-Artefakt} (z.\,B.\ CRP-Erh\"ohung durch chronische Entz\"undung, nicht akute Infektion), \textit{Kontraindikation} (z.\,B.\ Schwangerschaft bei bestimmter Medikation), \textit{zeitliche Fehlzuordnung} (z.\,B.\ Symptombeginn vor oder nach vermutetem Ausl\"oser), \textit{diagnostische Einschr\"ankung} (z.\,B.\ nicht verwertbarer Testbefund) und \textit{sonstiger Warnhinweis}. Jeder Warnhinweis verweist auf die betroffene Quellenachse (I--VI), sodass die klinische Relevanz sofort kontextualisiert ist. CAVE-Warnhinweise werden in der Gesamtsynopse prominent rot hervorgehoben.
\item[Longitudinaler Symptomverlauf] Eine strukturierte Dokumentation des zeitlichen Symptomverlaufs: F\"ur jedes relevante Symptom wird Beginn (Erstmanifestation), aktueller Status (aktiv/remittiert/fluktuierend) und Therapieansprechen (Ansprechen/kein Ansprechen/partielles Ansprechen) erfasst. Diese longitudinale Perspektive erm\"oglicht die Identifikation von Verlaufsmustern, Therapieresistenz und Chronifizierungsrisiken, die in einer rein querschnittlichen Diagnostik unsichtbar bleiben.
\item[Kontakt- und Beobachtungsprotokoll] Ein strukturiertes Protokoll f\"ur alle klinisch relevanten Kontakte und Beobachtungen. Jeder Eintrag umfasst Datum, Kontaktart (Telefonat, Gespr\"ach, Beobachtung, Hausbesuch, Fremdanamnese, E-Mail/Brief), beteiligte Kontaktperson, Inhalt und optionalen Achsenbezug. Das Protokoll dokumentiert den \textit{Prozess} der Informationsgewinnung: Wann wurde mit wem gesprochen, was wurde beobachtet, und welche Achse profitiert von dieser Information? Damit fungiert Achse~VI nicht nur als statische Belegsammlung, sondern als \textbf{Prozessdokumentation} des diagnostischen Arbeitens --- ein wesentlicher Bestandteil des Konzepts des lebenden Dokuments (vgl.\ Abschnitt~\ref{sec:designphilosophie}).
\end{description}

Die Erweiterung von Achse~VI reflektiert die klinische Erkenntnis, dass ein diagnostisches System nicht nur \textit{klassifizieren}, sondern auch \textit{warnen}, \textit{\"uber die Zeit verfolgen} und den \textit{diagnostischen Prozess selbst dokumentieren} muss. Die CAVE-Funktion adressiert das in der Polypharmazie- und Komplexit\"atsmedizin wachsende Problem klinischer Risiken, die in keiner einzelnen Achse vollst\"andig abgebildet sind, aber achsen\"ubergreifend Konsequenzen haben.

\subsection{Interprofessionelle Zuordnung}
\label{sec:multiprofessionell}

Das 6-Achsen-Modell bietet eine \textbf{empfohlene}, nicht verpflichtende Zuordnung von Achsen zu Professionsgruppen. Tabelle~\ref{tab:professionen} zeigt eine exemplarische Konfiguration f\"ur interdisziplin\"are Teams. In der klinischen Praxis k\"onnen die Zust\"andigkeiten flexibel angepasst werden: Ein Psychiater, der sowohl psychodiagnostische als auch somatische Kompetenz besitzt, kann selbstverst\"andlich Achse~I \textit{und} Achse~III ausf\"ullen; eine Fach\"arztin f\"ur Psychosomatische Medizin kann alle sechs Achsen bedienen. Das System \textit{erm\"oglicht} interdisziplin\"are Arbeitsteilung, \textit{erzwingt} sie aber nicht.

\begin{table}[H]
\centering
\caption{Exemplarische interprofessionelle Achsenzuordnung (nicht verpflichtend)}
\label{tab:professionen}
\begin{tabularx}{\textwidth}{lXX}
\toprule
\textbf{Achse} & \textbf{Typische Profession} & \textbf{Fachliche Begr\"undung} \\
\midrule
I & Psychologe / Psychiater & Psychodiagnostische Kompetenz; Testungshoheit \\
\addlinespace
II & Psychologe / Sozialp\"adagoge & Biographische Anamnese; Pers\"onlichkeitsdiagnostik \\
\addlinespace
III & Mediziner / Facharzt & Somatische Diagnosestellung; Kausalanalyse \\
\addlinespace
IV & Sozialarbeiter / Sozialp\"adagoge & Lebensweltexpertise; ICF-Kompetenz; Teilhabebewertung \\
\addlinespace
V & Interdisziplin\"ares Team & Synthese profitiert von allen Perspektiven \\
\addlinespace
VI & Alle Professionen & Jede beteiligte Profession belegt ihre eigenen Befunde \\
\bottomrule
\end{tabularx}
\end{table}

Der Mehrwert der Zuordnung liegt nicht in einer starren Zugangsbeschr\"ankung, sondern in der \textbf{strukturierten Aufforderung}: Das System macht transparent, \textit{welche} Kompetenz f\"ur welche Achse typischerweise ben\"otigt wird. In Einzelpraxen oder kleineren Settings kann eine einzelne Fachperson alle Achsen bearbeiten; in Kliniken und interdisziplin\"aren Teams erm\"oglicht das Modell eine effiziente Arbeitsteilung, bei der jede Profession in ihrer Dom\"ane mit identischen strukturellen Werkzeugen arbeitet. Die identische Subachsen-Struktur von Achse~I und~III ist dabei kein Zufall, sondern Designentscheidung: Ein Mediziner, der in IIIh eine somatische Verdachtsdiagnose formuliert, nutzt exakt dieselbe Logik wie ein Psychologe, der in Ih eine psychische Verdachtsdiagnose eintr\"agt.

\subsection{Komorbidit\"atsregeln und diagnostische Hierarchie}
\label{sec:komorbiditaet}

Das System implementiert drei Ebenen von Komorbidit\"atsregeln, die f\"ur die korrekte diagnostische Entscheidungsfindung essentiell sind:

\begin{description}[style=nextline, leftmargin=1.5cm, font=\bfseries]
\item[Hierarchische Exklusionsregeln] Bestimmte Diagnosen schlie\ss{}en andere aus. Beispiel: Eine Schizophrenie-Diagnose schlie\ss{}t eine gleichzeitige schizoaffektive St\"orung aus. Das System kodiert diese als harte Constraints, die bei der Diagnoseeingabe automatisch gepr\"uft werden.
\item[\glqq Due to another condition\grqq{}-Regeln] Viele DSM-5-TR-Diagnosen erfordern den Ausschluss, dass die Symptomatik durch eine andere psychische St\"orung, eine Substanz oder einen medizinischen Krankheitsfaktor besser erkl\"art wird. Diese Regeln werden direkt durch die Gatekeeper-Stufen~1--3 (Abschnitt~\ref{sec:gatekeeper}) implementiert.
\item[Erlaubte Komorbidit\"aten mit Priorisierung] Bei 4+ aktiven Diagnosen priorisiert das System nach: (1)~klinischer Dringlichkeit (Suizidalit\"at, Psychose), (2)~Schweregrad (Achse-I-Akutstatus), (3)~Behandelbarkeit (modifizierbare aufrechterhaltende Faktoren aus Achse~V), (4)~chronologischer Reihenfolge.
\end{description}


\clearpage
\subsection{Vollst\"andiges Sub-Achsen-Inventar}

Die Tabellen~\ref{tab:subaxes_part1}, \ref{tab:subaxes_part2} und \ref{tab:subaxes_part3} bieten ein vollst\"andiges Inventar aller Sub-Achsen der sechs Hauptachsen einschlie\ss{}lich der zugeordneten Instrumente und Datenquellen.

\begin{table}[H]
\centering
\caption{Sub-Achsen-Inventar: Achsen I--II}
\label{tab:subaxes_part1}
\begin{CompactLongTable}
\begin{tabular}{@{}L{1.8cm}L{2.4cm}L{5.0cm}L{4.0cm}@{}}
\toprule
\textbf{Achse} & \textbf{Sub-Achse} & \textbf{Beschreibung} & \textbf{Instrument / Quelle} \\
\midrule
I\\Psychisch & Ia: Akute St\"orungen & Aktuelle DSM-5-TR/ICD-11-Diagnosen mit PRO/CONTRA-Evidenz und Konfidenz & DSM-5-TR, ICD-11 \\
 & Ib: Chronische Verl\"aufe & Langzeitdiagnosen mit Chronifizierungsrisiko & Klinische Akten \\
 & Ic: Widerlegte Verdachte & Explizit ausgeschlossene Diagnosen mit differenzialdiagn. Begr\"undung & Differenzialdiagnostik \\
 & Id: Remittierte Diagnosen & Fr\"uhere Diagnosen mit belegter Inaktivit\"at & Klinische Akten \\
 & Ie: Remissionsfaktoren & Therapie, Medikation oder spontane Bew\"altigung als Remissionsursache & Klinische Akten \\
 & If: Behandlungsgeschichte & Fr\"uhere Behandlungen, Medikation, Wirkungen und Nebenwirkungen & Klinische Akten \\
 & Ig: Therapietreue & Adh\"arenz aus Selbst- und Fremdperspektive & Klinisches Interview \\
 & Ih: Verdachtsdiagnosen & Laufende diagnostische Hypothesen (noch nicht best\"atigt) & Klinisches Urteil \\
 & Ii: Abdeckungsanalyse & Symptom-Diagnosen-Abdeckungsprozentsatz (DCI) & Algorithmisch (s.\ Abschnitt~\ref{sec:abdeckung}) \\
 & Ij: Untersuchungsplan & Priorisierte n\"achste Schritte (dringend / wichtig / Monitoring) & Klinisches Urteil \\
\midrule
II\\Pers\"onlichkeit & IIa: PID-5 Dimensionales Profil & 5 Trait-Dom\"anen, 25 Facetten & PID-5, PID5BF+ \\
 & IIb: Pr\"agende Lebenserfahrungen & Biographische Ereignisse mit Einfluss auf Pers\"onlichkeitsentwicklung & Klinisches Interview \\
 & IIc: Grundkonflikte & Zentrale Beziehungsmuster und Konfliktthemen & OPD-3, klin. Formulierung \\
\bottomrule
\end{tabular}
\end{CompactLongTable}
\end{table}

\begin{table}[H]
\centering
\caption{Sub-Achsen-Inventar: Achse III}
\label{tab:subaxes_part2}
\begin{CompactLongTable}
\begin{tabular}{@{}L{1.8cm}L{2.4cm}L{5.0cm}L{4.0cm}@{}}
\toprule
\textbf{Achse} & \textbf{Sub-Achse} & \textbf{Beschreibung} & \textbf{Instrument / Quelle} \\
\midrule
III\\Somatisch & IIIa: Akute somatische Diagnosen & Aktuelle somatische Diagnosen mit PRO/CONTRA-Evidenz & ICD-11 somatische Codes \\
 & IIIb: Chronische Diagnosen (erkl\"arend) & Erkrankungen, die bestimmte Symptome vollst\"andig erkl\"aren & Klinische Akten \\
 & IIIc: Beitragende med. Faktoren & Somatische Befunde, die Psychopathologie beg\"unstigen oder verst\"arken & Klinische Akten \\
 & IIId: Remittierte Erkrankungen & Fr\"uhere somatische Erkrankungen mit belegter Inaktivit\"at & Klinische Akten \\
 & IIIe: Remissionsfaktoren & Operative Eingriffe, Medikation, Lebensstil\"anderung als Remissionsursache & Klinische Akten \\
 & IIIf: Behandlungsgeschichte & Fr\"uhere medizinische Behandlungen, Operationen, Nebenwirkungen & Klinische Akten \\
 & IIIg: Medikamentenadh\"arenz & Adh\"arenz aus \"arztlicher Dokumentation und Patientenbericht & Klinische Akten \\
 & IIIh: Verdachtsdiagnosen & Laufende medizinische diagnostische Hypothesen & Klinisches Urteil \\
 & IIIi: Abdeckungsanalyse (somatisch) & Somatische Symptom-Diagnosen-Abdeckung & Algorithmisch (s.\ Abschnitt~\ref{sec:abdeckung}) \\
 & IIIj: Somatischer Untersuchungsplan & Medizinische Abkl\"arungspriorit\"aten & Klinisches Urteil \\
 & IIIk: Integrierte Kausalanalyse & Achsen\"ubergreifende Kausalpfade (somatisch $\leftrightarrow$ psychisch) & Interdisziplin\"are Synthese \\
 & IIIl: Genetik/Familie & Heredit\"are Risikofaktoren und Familienanamnese & Familien\-anamnese-Interview \\
 & IIIm: Medikamentenanamnese & Aktuelle und fr\"uhere Medikation mit Wirksamkeitsbewertung & Medikamentenakten \\
\bottomrule
\end{tabular}
\end{CompactLongTable}
\end{table}

\begin{table}[H]
\centering
\caption{Sub-Achsen-Inventar: Achsen IV--VI}
\label{tab:subaxes_part3}
\begin{CompactLongTable}
\begin{tabular}{@{}L{1.8cm}L{2.4cm}L{5.0cm}L{4.0cm}@{}}
\toprule
\textbf{Achse} & \textbf{Sub-Achse} & \textbf{Beschreibung} & \textbf{Instrument / Quelle} \\
\midrule
IV\\Funktions-\\status & IVa: WHODAS 2.0 & 6 Funktionsdom\"anen (36- oder 12-Item-Version) & WHODAS 2.0 \\
 & IVb: GdB & Grad der Behinderung & VersMedV-Richtlinien \\
 & IVc: ICF Core Sets & St\"orungsspezifische ICF-Komponentenprofile & ICF Core Sets \\
 & IVd: Cultural Formulation Interview & Kultureller Kontext der Krankheitserfahrung & CFI (DSM-5-TR) \\
 & IVe: Behandlungsnetzwerk\\Register & Kontaktpersonen, Versorger, soziale Unterst\"utzung & Strukturiertes Register \\
\midrule
V\\Fall-\\formulierung & Va: Pr\"adisponierende Faktoren & Vulnerabilit\"atsfaktoren vor St\"orungsbeginn & 3P/4P-Modell \\
 & Vb: Ausl\"osende Faktoren & Triggerereignisse und Stressoren & 3P/4P-Modell \\
 & Vc: Aufrechterhaltende Faktoren & Aufrechterhaltende Mechanismen und R\"uckkopplungsschleifen & 3P/4P-Modell \\
 & Vd: Protektive Faktoren & Resilienzressourcen und St\"arken & 3P/4P-Modell \\
 & Ve: Abdeckungsl\"ucken-Integration & Unerkl\"arte Symptome aus Achse-I/III-Abdeckungsanalyse & Achsen\"ubergreifende Synthese \\
\midrule
VI\\Evidenz\\\& Sicherheit & VIa: Zentrale Evidenz-Matrix & Strukturierte Evidenzfelder pro Diagnose & Systemgeneriert \\
 & VIb: CAVE Klinische Warnungen & Achsen\"ubergreifende Risikowarnungen (Interaktionen, Kontraindikationen) & Regelbasierte Warnhinweise \\
 & VIc: Longitudinale Symptomverlaufsdokumentation & Wiederholte Messungen mit Therapieansprechen-Dokumentation & PHQ-9, GAD-7, PCL-5 etc. \\
 & VId: Kontakt- und Beobachtungsprotokoll & Sitzungsdokumentation und Prozessnotizen & Kliniker-Eintr\"age \\
\bottomrule
\end{tabular}
\end{CompactLongTable}
\end{table}


% ===================================================================
% 3. GATEKEEPER-LOGIK
% ===================================================================
\section{Die 6-Stufen-Gatekeeper-Logik der Differenzialdiagnostik}
\label{sec:gatekeeper}

Das System implementiert eine dynamische Warteschlange, in der Screening-Auff\"alligkeiten spezifische DSM-5-TR/ICD-11-Module ausl\"osen. Die 6-Stufen-Sequenz bildet exakt die von Michael~B. First im \textit{DSM-5-TR Handbook of Differential Diagnosis} (2024) publizierte differenzialdiagnostische Sequenz ab (Tabelle~\ref{tab:gatekeeper}).

\begin{table}[H]
\centering
\caption{Abgleich der 6-Stufen-Gatekeeper-Logik mit Firsts Goldstandard-Sequenz}
\label{tab:gatekeeper}
\begin{tabularx}{\textwidth}{clXc}
\toprule
\textbf{Stufe} & \textbf{Systemschritt} & \textbf{Firsts Sequenz} & \textbf{Passung} \\
\midrule
1 & Simulationsausschluss & Malingering und artifizielle St\"orung ausschlie\ss{}en & Exakt \\
\addlinespace
2 & Substanzausschluss & Substanz\"atiologie ausschlie\ss{}en & Exakt \\
\addlinespace
3 & Medizinischer Ausschluss & \"Atiologische medizinische Erkrankung ausschlie\ss{}en & Exakt \\
\addlinespace
4 & Prim\"arkategorie & Spezifische Prim\"arst\"orung(en) bestimmen & Exakt \\
\addlinespace
5 & Anpassungsst\"orung & Anpassungsst\"orungen differenzieren & Exakt \\
\addlinespace
6 & Funktionsschwelle & Grenze zu \glqq keine psychische St\"orung\grqq{} & Konzeptuell \\
\bottomrule
\end{tabularx}
\end{table}

\citet{First2024} beschreibt diese Sequenz als den definitiven Makro-Level-Diagnoseprozess. Sein Handbook enth\"alt dar\"uber hinaus \textbf{30 symptomorientierte Entscheidungsb\"aume} (zwei in der TR-Edition erg\"anzt f\"ur dissoziative Symptome und repetitive pathologische Verhaltensweisen) und \textbf{67~Differenzialdiagnosetabellen}, die alle dieser Sequenz folgen.

\textbf{Anmerkung zu Stufe~6 (\glqq konzeptuell\grqq{}):} W\"ahrend die Stufen~1--5 exakt algorithmisch abbildbar sind, ist Stufe~6 --- die Grenzziehung zwischen normaler Stressreaktion und psychischer St\"orung --- inh\"arent eine klinische Urteilsentscheidung. Das System unterst\"utzt diese durch objektive Funktionsdaten (WHODAS~2.0, Mini-ICF-APP aus Achse~IV) und die Cross-Cutting-Schwellenwerte, kann aber die klinische Entscheidung nicht vollst\"andig automatisieren. Diese epistemische Grenze wird transparent kommuniziert.

\begin{figure}[ht]
\centering
\begin{tikzpicture}[
    stepbox/.style={draw, rounded corners=3pt, minimum width=5.5cm, minimum height=0.8cm,
        font=\small, align=center, fill=blue!8},
    decision/.style={diamond, draw, aspect=2.2, minimum width=2.2cm,
        font=\scriptsize, align=center, fill=yellow!15, inner sep=1pt},
    exitbox/.style={draw, rounded corners=3pt, minimum width=3cm, minimum height=0.6cm,
        font=\scriptsize, align=center, fill=red!10},
    conceptual/.style={draw, rounded corners=3pt, dashed, minimum width=5.5cm, minimum height=0.8cm,
        font=\small, align=center, fill=gray!10},
    >=Stealth, node distance=1.2cm
]

% Stufen
\node[stepbox] (s1) {Stufe 1: Ausschluss\\Simulation / Artifiziell};
\node[decision, right=1.5cm of s1] (d1) {Ja?};
\node[exitbox, right=1.2cm of d1] (e1) {Simulation /\\Artif. St\"orung};

\node[stepbox, below=1.0cm of s1] (s2) {Stufe 2: Ausschluss\\Substanz\"atiologie};
\node[decision, right=1.5cm of s2] (d2) {Ja?};
\node[exitbox, right=1.2cm of d2] (e2) {Substanz-\\induzierte Dx};

\node[stepbox, below=1.0cm of s2] (s3) {Stufe 3: Ausschluss\\Med. Erkrankung};
\node[decision, right=1.5cm of s3] (d3) {Ja?};
\node[exitbox, right=1.2cm of d3] (e3) {Med.\\Erkrankung};

\node[stepbox, below=1.0cm of s3] (s4) {Stufe 4: Bestimmung\\Prim\"arst\"orung(en)};
\node[decision, right=1.5cm of s4] (d4) {Ja?};
\node[exitbox, right=1.2cm of d4] (e4) {Prim\"are\\St\"orung(en)};

\node[stepbox, below=1.0cm of s4] (s5) {Stufe 5: Differenzierung\\Anpassungsst\"orung};
\node[decision, right=1.5cm of s5] (d5) {Ja?};
\node[exitbox, right=1.2cm of d5] (e5) {Anpassungs-\\st\"orung};

\node[conceptual, below=1.0cm of s5] (s6) {Stufe 6: Grenze zu\\keiner psych. St\"orung \textnormal{\scriptsize(konzeptuell)}};

% Pfeile: Stufe zu Entscheidung
\draw[->] (s1) -- (d1);
\draw[->] (s2) -- (d2);
\draw[->] (s3) -- (d3);
\draw[->] (s4) -- (d4);
\draw[->] (s5) -- (d5);

% Pfeile: Entscheidung zu Ausgang (Ja)
\draw[->] (d1) -- node[above, font=\scriptsize] {Ja} (e1);
\draw[->] (d2) -- node[above, font=\scriptsize] {Ja} (e2);
\draw[->] (d3) -- node[above, font=\scriptsize] {Ja} (e3);
\draw[->] (d4) -- node[above, font=\scriptsize] {Ja} (e4);
\draw[->] (d5) -- node[above, font=\scriptsize] {Ja} (e5);

% Pfeile: Entscheidung Nein -> n\"achste Stufe
\draw[->] (d1.south) -- ++(0,-0.3) -| node[near start, right, font=\scriptsize] {Nein} (s2.north);
\draw[->] (d2.south) -- ++(0,-0.3) -| node[near start, right, font=\scriptsize] {Nein} (s3.north);
\draw[->] (d3.south) -- ++(0,-0.3) -| node[near start, right, font=\scriptsize] {Nein} (s4.north);
\draw[->] (d4.south) -- ++(0,-0.3) -| node[near start, right, font=\scriptsize] {Nein} (s5.north);
\draw[->] (d5.south) -- ++(0,-0.3) -| node[near start, right, font=\scriptsize] {Nein} (s6.north);

\end{tikzpicture}
\caption{6-Stufen-Gatekeeper-Logik f\"ur die Differenzialdiagnostik, basierend auf \citet{First2024}. Stufe~6 ist als konzeptuell markiert, da die Grenzziehung zwischen normaler Variation und psychischer St\"orung klinisches Urteil erfordert.}
\label{fig:gatekeeper}
\end{figure}


% ===================================================================
% 4. CROSS-CUTTING SCREENING
% ===================================================================
\section{Cross-Cutting Symptom Measures als intelligente Triage}
\label{sec:screening}

Die DSM-5 Cross-Cutting Symptom Measures bilden das R\"uckgrat der Screening-Architektur. Das \textbf{Level-1-Erwachsenenma\ss{}} enth\"alt \textbf{23~Items \"uber 13~Dom\"anen}: Depression (2), \"Arger (1), Manie (2), Angst (3), somatische Symptome (2), Suizidalit\"at (1), Psychose (2), Schlafprobleme (1), Ged\"achtnis (1), repetitive Gedanken und Verhaltensweisen (2), Dissoziation (1), Pers\"onlichkeitsfunktion (2) und Substanzgebrauch (3). Jedes Item verwendet eine 5-stufige Likert-Skala (0~=~keine bis 4~=~schwer).

Die Schwellenlogik ist klinisch kalibriert: \textbf{Die meisten Dom\"anen l\"osen Level~2 bei $\geq 2$ (leicht) aus}, aber drei sicherheitskritische Dom\"anen --- Suizidalit\"at, Psychose und Substanzgebrauch --- \textbf{l\"osen bei $\geq 1$ (gering) aus}. Diese Asymmetrie reflektiert den klinischen Imperativ, selbst minimale Best\"atigung gef\"ahrlicher Symptome nicht zu \"ubersehen.

Die Level-2-Ma\ss{}e verweisen auf spezifische validierte Instrumente: PROMIS Depression Short Form, PROMIS Anxiety Short Form, Altman Self-Rating Mania Scale (ASRM), PHQ-15 f\"ur somatische Symptome, PROMIS Sleep Disturbance, adaptierte FOCI Severity Scale f\"ur Zwangsst\"orungen und adaptierter NIDA-Modified ASSIST f\"ur Substanzgebrauch. F\"unf Dom\"anen (Suizidalit\"at, Psychose, Ged\"achtnis, Dissoziation, Pers\"onlichkeitsfunktion) besitzen \textit{keine offiziellen Level-2-Ma\ss{}e} --- die APA empfiehlt hier klinische Evaluation.

Die DSM-5 Field Trials \citep{Narrow2013} zeigten \textbf{gute bis exzellente Test-Retest-Reliabilit\"at} (ICC~0,64--0,97) f\"ur die meisten Items.

\subsection{Kinder- und Jugendversion}
\label{sec:screening_kj}

Das DSM-5 Cross-Cutting-System umfasst eine separate \textbf{Elternberichtsversion f\"ur Kinder und Jugendliche} (6--17~Jahre) mit \textbf{25~Items \"uber 12~Dom\"anen}. Die Dom\"ane \glqq Pers\"onlichkeitsfunktion\grqq{} entf\"allt entwicklungsbedingt; stattdessen werden Reizbarkeit und \"Arger st\"arker gewichtet. Eine vollst\"andige Implementierung erfordert die p\"adiatrische Variante als separaten Einstiegspfad sowie altersangepasste Level-2-Instrumente.

Das vorliegende System bevorzugt \textbf{ausschlie\ss{}lich frei verf\"ugbare Instrumente}; der propriet\"are CBCL/YSR-Komplex \citep{Achenbach2001} wird zwar als Referenzstandard anerkannt, aber durch lizenzfreie Alternativen ersetzt. Tabelle~\ref{tab:instrumente_kj} listet die empfohlene Auswahl freier K/J-Instrumente parallel zur Erwachsenenmatrix (Tabelle~\ref{tab:instrumente}).

\begin{table}[htbp]
\centering
\caption{Freie Screening-Instrumente f\"ur Kinder und Jugendliche (K/J-Matrix)}
\label{tab:instrumente_kj}
\begin{CompactLongTable}
\begin{tabular}{@{}L{2.4cm}L{2.4cm}L{0.9cm}L{1.5cm}L{2.0cm}L{1.5cm}@{}}
\toprule
\textbf{Dom\"ane} & \textbf{Instrument} & \textbf{Items} & \textbf{Alter} & \textbf{Cutoff / Norm} & \textbf{Kosten} \\
\midrule
Cross-Cutting (Eltern) & DSM-5 XC L1 Parent & 25 & 6--17 & $\geq 1$/$\geq 2$ Triage & Frei \\
\addlinespace
Allg. psychosoz. Screen & SDQ (Self/Parent) & 25 & 4--17 & $\geq 20$ (Total) & Frei \\
\addlinespace
Depression & PHQ-9 mod. (PHQ-A) & 9 & 11--17 & $\geq 11$ & Frei \\
\addlinespace
Angst + Depression & RCADS-25 & 25 & 8--18 & $T \geq 65$ & Frei \\
\addlinespace
Angst (multidim.) & SCARED-C & 41 & 8--18 & $\geq 25$ & Frei \\
\addlinespace
PTBS & CPSS-V / CRIES-8 & 20 / 8 & 7--18 & $\geq 31$ / $\geq 17$ & Frei \\
\addlinespace
Substanzgebrauch & CRAFFT 2.1+N & 6+3 & 12--21 & $\geq 2$ & Frei \\
\addlinespace
ADHS & SNAP-IV-26 & 26 & 6--18 & Subskalen-Mittelwert & Frei \\
\addlinespace
Autismus & AQ-10 Adolescent/Child & 10 & 4--15 & $\geq 6$ & Frei \\
\addlinespace
Zwang & CY-BOCS-SR & 10 & 8--18 & $\geq 16$ & Frei \\
\addlinespace
Suizidalit\"at & C-SSRS Pediatric & 6 & $\geq 5$ & Jede Best\"at. & Frei \\
\addlinespace
Essst\"orungen & SCOFF (ab $\sim$13~J.) & 5 & 13--18 & $\geq 2$ & Frei \\
\bottomrule
\end{tabular}
\end{CompactLongTable}
\end{table}

Diese Auswahl bewahrt die diagnostische Logik der Erwachsenenmatrix (Dom\"anen-orientiert, Cutoff-basiert), substituiert aber jedes Instrument durch eine entwicklungspsychologisch validierte und lizenzfreie K/J-Variante. F\"ur sicherheitskritische Dom\"anen (Suizidalit\"at, Substanzgebrauch, Psychose) gilt analog die niedrigere Triage-Schwelle. Die Trennung in \textbf{Selbstbericht} (ab $\sim$8--11~J.) und \textbf{Elternbericht} (6--17~J., bei jungen Kindern obligatorisch) folgt der etablierten Konvention der p\"adiatrischen Psychometrie.

\subsection{Umfassende Screening-Instrument-Matrix}

Tabelle~\ref{tab:instrumente} zeigt die empfohlenen Instrumente f\"ur alle wesentlichen diagnostischen Dom\"anen.

\begin{table}[htbp]
\centering
\caption{Screening-Instrument-Matrix}
\label{tab:instrumente}
\begin{CompactLongTable}
\begin{tabular}{@{}L{2.4cm}L{2.0cm}L{0.9cm}L{1.8cm}L{2.2cm}L{1.1cm}@{}}
\toprule
\textbf{Dom\"ane} & \textbf{Instrument} & \textbf{Items} & \textbf{Cutoff} & \textbf{Sens./Spez.} & \textbf{Kosten} \\
\midrule
Depression & PHQ-9 & 9 & $\geq 10$ & 88\%/88\% & Frei \\
\addlinespace
Angst & GAD-7 & 7 & $\geq 10$ & 89\%/82\% & Frei \\
\addlinespace
PTBS (DSM-5) & PCL-5 & 20 & 31--33 & 85--95\%/82--90\% & Frei \\
\addlinespace
PTBS/KPTBS (ICD-11) & ITQ & 18 & Algorithmus & Exzellent & Frei \\
\addlinespace
Psychoserisiko & PQ-16 & 16 & $\geq 6$ & 87\%/87\% & Frei \\
\addlinespace
Bipolar-Screening & MDQ & 15 & $\geq 7$ & 73\%/90\% & Frei \\
\addlinespace
ADHS & ASRS~v1.1 & 6 & $\geq 4$ & 69\%/99,5\% & Frei \\
\addlinespace
Autismus & AQ-10 & 10 & $\geq 6$ & 88\%/91\% & Frei \\
\addlinespace
Alkoholgebrauch & AUDIT & 10 & $\geq 8$ & 92\%/94\% & Frei \\
\addlinespace
Drogengebrauch & DAST-10 & 10 & $\geq 3$ & 98\%/91\% & Frei \\
\addlinespace
Pers\"onlichkeit & PID-5-BF & 25 & Dimensional & N/A & Frei \\
\addlinespace
Suizidalit\"at & C-SSRS & 6 & Jede Best\"at. & Risikoklassif. & Frei \\
\addlinespace
Zwang & OCI-R & 18 & $\geq 21$ & Gut & Frei \\
\addlinespace
Somatisierung & SSS-8 & 8 & $\geq 12$ & Vergleichbar & Frei \\
\addlinespace
Dissoziation & DES-II & 28 & $\geq 30$ & 74\%/80\% & Frei \\
\addlinespace
Essst\"orungen & SCOFF & 5 & $\geq 2$ & 84\%/90\% & Frei \\
\addlinespace
Essst\"orungen (detail.) & EDE-QS & 12 & $\geq 15$ & Gut/Gut & Frei \\
\addlinespace
Schlafst\"orungen & ISI & 7 & $\geq 15$ & 82\%/82\% & Frei \\
\addlinespace
Funktionsf\"ahigkeit & WHODAS 2.0 (12) & 12 & Dimensional & N/A & Frei \\
\bottomrule
\end{tabular}
\end{CompactLongTable}
\end{table}

Alle empfohlenen Instrumente sind \textbf{frei verf\"ugbar} --- keine Lizenzkosten behindern die Implementierung. Die Auswahl basiert auf den jeweiligen Validierungspublikationen: PHQ-9 \citep{Levis2019}, GAD-7 \citep{Spitzer2006}, PCL-5 \citep{Blevins2015}, ITQ \citep{Cloitre2018}, PQ-16 \citep{Ising2012}, ASRS \citep{Kessler2005}, AQ-10 \citep{Allison2012}, AUDIT \citep{Saunders1993}, PID-5 \citep{OltmannsWidiger2018, Riegel2021}, C-SSRS \citep{Posner2011}, OCI-R \citep{Foa2002}, SSS-8 \citep{Gierk2014}, DES-II \citep{CarlsonPutnam1993}, SCOFF \citep{Morgan1999}, ISI \citep{Morin2011}. F\"ur die Interrater-Reliabilit\"at kategorialer DSM-5-Diagnosen siehe \citet{Regier2013} und \citet{Aboraya2006}.


% ===================================================================
% 5. KLASSIFIKATIONSDIVERGENZEN
% ===================================================================
\section{Kritische Divergenzen zwischen DSM-5-TR und ICD-11}
\label{sec:divergenzen}

Das System muss f\"unf kritische Divergenzen zwischen den Klassifikationssystemen kodieren \citep{Keeley2016}.

\subsection{Schizophrenie-Dauer}

DSM-5-TR (295.90/F20.x) verlangt \textbf{6~Monate} kontinuierlicher Zeichen einschlie\ss{}lich mindestens 1~Monat aktiver Symptome, w\"ahrend ICD-11 (6A20) nur \textbf{1~Monat} fordert. Ein Patient kann somit die ICD-11-Kriterien f\"ur Schizophrenie erf\"ullen, unter DSM-5 aber nur die Diagnose einer schizophreniformen St\"orung erhalten.

\subsection{PTBS-Architektur}

DSM-5-TR verwendet \textbf{4~Cluster mit 20~Symptomen} (Intrusion $\geq$1/5, Vermeidung $\geq$1/2, Negative Kognitionen $\geq$2/7, Arousal $\geq$2/6). ICD-11 verwendet ein bewusst schmaleres Modell: \textbf{3~Cluster mit 6~Kernsymptomen}. ICD-11 f\"ugt dann die \textbf{Komplexe PTBS (6B41)} als distinkte Diagnose hinzu, die alle PTBS-Kriterien plus drei St\"orungen der Selbstorganisation (Affektdysregulation, negatives Selbstkonzept, Beziehungsschwierigkeiten) erfordert.

\subsection{Pers\"onlichkeitsst\"orungen}

ICD-11 ersetzte kategoriale Typen vollst\"andig durch ein dimensionales Modell. Das DSM-5 Alternative Model for Personality Disorders (AMPD) verwendet die LPFS f\"ur Kriterium~A (Funktionsniveau) und f\"unf Trait-Dom\"anen mit 25~Facetten f\"ur Kriterium~B (pathologische Traits). Der PID-5-BF+M verbr\"uckt beide Systeme. \textbf{Abgrenzung zu Achse~II des vorliegenden Modells:} Das AMPD selbst ist implizit zweiachsig (Funktionsniveau + Traits). Achse~II des 6-Achsen-Modells \"ubernimmt die Trait-Dimension (PID-5), erweitert sie aber um biographische Kontextualisierung (pr\"agende Erfahrungen, Grundkonflikte) und verlagert das Funktionsniveau (LPFS-Analogon) systematisch auf Achse~IV (WHODAS~2.0, Mini-ICF-APP). Dadurch wird die diagnostische Information des AMPD vollst\"andig abgebildet, aber \"uber mehrere Achsen verteilt, was eine klarere Zuordnung zu unterschiedlichen professionellen Zust\"andigkeiten erm\"oglicht.

\subsection{Gaming Disorder}

ICD-11 (6C51) verwendet einen monothetischen Ansatz (alle 4~Kriterien erforderlich; $\geq 12$~Monate). DSM-5-TR listet Internet Gaming Disorder nur als \glqq Condition for Further Study\grqq{} mit polythetischem Ansatz ($\geq 5$ von 9~Kriterien). Konkordanz: $\kappa = 0{,}80$, aber ICD-11 hat eine h\"ohere diagnostische Schwelle (Pr\"avalenz 2,7\% vs.\ DSM-5 5,2\%).

\subsection{Komplexe Trauer / Prolonged Grief Disorder}

ICD-11 (6B42) und DSM-5-TR (Prolonged Grief Disorder, neu aufgenommen) kodieren Anh\"altende Trauerst\"orung, unterscheiden sich aber in Zeitkriterium (ICD-11: 6~Monate; DSM-5-TR: 12~Monate) und Symptomkonfiguration. Das System muss beide Varianten abbilden.


% ===================================================================
% 6. ABDECKUNGSANALYSE
% ===================================================================
\section{Die Abdeckungsanalyse: Ein vorgeschlagener Ansatz}
\label{sec:abdeckung}

Die Abdeckungsanalyse (\textit{Coverage Analysis}) stellt eine zentrale vorgeschlagene Komponente des Systems dar. Uns sind keine direkt vergleichbaren Implementierungen bekannt, obwohl verwandte Ans\"atze existieren (z.\,B.\ Bayesianische Symptom-Diagnosen-Zuordnung in OPCRIT). Das Konzept umfasst die systematische Kennzeichnung von Symptomen, die durch aktuelle Diagnosen nicht erkl\"art werden.

Die engsten existierenden Parallelen sind: Komorbidit\"ats-Detektionsalgorithmen, die Symptome kennzeichnen, die auf zus\"atzliche Diagnosen hindeuten; Firsts Differenzialdiagnosetabellen \citep{First2024}, die \"uberlappende Pr\"asentationen vergleichen; und transdiagnostische Rahmenmodelle wie HiTOP, in denen Symptome als Netzwerkknoten modelliert werden. \citet{PlanaRipoll2019} zeigten empirisch, dass die meisten psychischen St\"orungen weit h\"aufiger gemeinsam auftreten als durch Zufall erwartet --- \textbf{eine vorangegangene Diagnose irgendeiner psychischen St\"orung erh\"ohte das Risiko f\"ur alle anderen psychischen St\"orungen substanziell} --- was den Bedarf an einem systematischen Abdeckungsanalyse-Werkzeug direkt st\"utzt.

\subsection{Formale Definition der Abdeckungsanalyse}

Sei $S = \{s_1, \ldots, s_n\}$ die Menge der beobachteten Symptome und $D = \{d_1, \ldots, d_m\}$ die Menge der etablierten Diagnosen. F\"ur jedes Symptom $s_i$ definieren wir die Abdeckungsfunktion $c(s_i) = \max_{d_j \in D} w(s_i, d_j)$, wobei $w(s_i, d_j) \in [0, 1]$ das Gewicht des Symptoms $s_i$ in der Diagnose $d_j$ gem\"a\ss{} den diagnostischen Kriterien repr\"asentiert. Die Gesamtabdeckung --- im Folgenden als \textbf{Diagnostic Coverage Index (DCI)} bezeichnet --- ergibt sich als $C_{\text{total}} = \frac{1}{n} \sum_{i=1}^{n} c(s_i)$.

Die Symptomgewichte werden aus den diagnostischen Kriterien von DSM-5-TR und ICD-11 abgeleitet. Kriterien werden nach diagnostischer Salienz gewichtet: Kardinalsymptome erhalten hohe Gewichte ($w \geq 0{,}8$), w\"ahrend unspezifische Symptome niedrigere Gewichte ($0{,}3 \leq w \leq 0{,}7$) erhalten. Dieses initiale Gewichtungsschema ist heuristisch und bedarf empirischer Kalibrierung in zuk\"unftiger Arbeit.

\"Uberlappende Diagnosen werden behandelt, indem jedes Symptom der Diagnose mit maximaler Abdeckung zugeordnet wird. Die Wahl der Max-Funktion (statt z.\,B.\ eines gewichteten Mittels oder einer Bayesianischen Posterior-Wahrscheinlichkeit) ist bewusst konservativ: Sie fragt, ob \textit{mindestens eine} Diagnose das Symptom hinreichend erkl\"art, und vermeidet dadurch die Summation partieller Erkl\"arungen, die zu falsch hohen Abdeckungswerten f\"uhren k\"onnte. Dieser Ansatz reflektiert die klinische Praxis, in der ein Symptom als \glqq erkl\"art\grqq{} gilt, wenn es einem etablierten Krankheitsbild zugeordnet werden kann. Alternative Aggregationsfunktionen (gewichtete Summe, Bayesianische Posterior) sollten in zuk\"unftigen Studien systematisch verglichen werden.

\subsection{Begr\"undung der Schwellenwerte und Sensitivit\"at}

Die vorgeschlagenen Schwellenwerte (85\%/60\%) sind als \textit{vorl\"aufige, heuristische Ausgangswerte} zu verstehen, die auf Analogien zu verwandten, aber methodisch distinkten Metriken basieren: (1)~In der somatischen Medizin gelten diagnostische Sensitivit\"aten $\geq$85\% als klinisch akzeptabel f\"ur Screening-Zwecke \citep{Levis2019}; dieser Schwellenwert wurde auf die diagnostische Abdeckung \"ubertragen. (2)~Der untere Schwellenwert von 60\% orientiert sich an der Konvention, dass Fleiss'~$\kappa$-Werte $<$0,6 als ``unzureichende \"Ubereinstimmung'' klassifiziert werden. Es ist ausdr\"ucklich anzumerken, dass diagnostische Abdeckung ein anderes Konstrukt ist als Screening-Sensitivit\"at oder Interrater-\"Ubereinstimmung; die Analogien dienen lediglich als orientierende Anhaltspunkte, nicht als methodische Begr\"undung. Diese Ausgangswerte m\"ussen zwingend durch empirische Kalibrierung in mindestens drei klinischen Kohorten (station\"ar, ambulant, somatisch komorbid) ersetzt werden. Hinsichtlich der Sensitivit\"at der Gesamtmetrik: Da $C_{\text{total}}$ ein arithmetisches Mittel \"uber $n$ Symptome ist, ver\"andert eine Einzelgewichts\"anderung von $\Delta w$ den Gesamtwert um $\Delta C = \Delta w / n$. Bei typischen Symptomzahlen ($n = 8$--$15$) f\"uhrt eine Fehlsch\"atzung eines einzelnen Gewichts um $\pm 0{,}2$ zu einer Gesamtverschiebung von lediglich $\pm 1{,}3$--$2{,}5$ Prozentpunkten --- was die Robustheit der Metrik gegen\"uber moderaten Gewichtungsfehlern belegt, aber auch zeigt, dass bei wenigen Symptomen ($n < 5$) die Metrik instabil wird.

\subsection{Implementierung}

Die Implementierung arbeitet wie folgt: (1)~Sammlung aller best\"atigten Symptome \"uber alle Screening-Instrumente, (2)~Zuordnung jedes Symptoms zu den diagnostischen Kriterien, die es erf\"ullt, (3)~Berechnung eines Abdeckungsprozentsatzes pro Symptom, der den Grad der diagnostischen Erkl\"arung quantifiziert, (4)~Klassifikation als vollst\"andig abgedeckt ($\geq$85\%), partiell abgedeckt (60--84\%) oder unzureichend abgedeckt ($<$60\%), (5)~Berechnung eines Gesamtabdeckungswerts als gewichteter Mittelwert \"uber alle Symptome, (6)~Pr\"asentation der Ergebnisse als Symptom-Diagnosen-Matrix mit visueller Hervorhebung diagnostischer L\"ucken. Das 4P-Modell (Achse~V) dient als nat\"urliches Komplement: Symptome, die nicht durch Achse-I-Diagnosen abgedeckt sind (Achse~Ii), sollten durch die Achse-V-Formulierung erkl\"arbar sein; solche, die es nicht sind, repr\"asentieren genuine diagnostische L\"ucken und m\"unden in den priorisierten Untersuchungsplan (Achse~Ij).

Die quantitative Metrik (z.\,B.\ \glqq Gesamtabdeckung: $\sim$88\%\grqq{}) liefert dem Kliniker eine unmittelbare R\"uckmeldung \"uber die Vollst\"andigkeit seiner diagnostischen Arbeit. Verwandte Ans\"atze existieren in Bayesianischen Diagnosesystemen (z.\,B.\ DXplain, Isabel Healthcare), die Symptom-Diagnose-Zuordnungen vornehmen; der spezifische Beitrag der vorliegenden Abdeckungsanalyse liegt in der transparenten, prozentbasierten Darstellung als klinische Qualit\"atssicherungsmetrik.

Es sei angemerkt, dass der Begriff \glqq Abdeckungsanalyse\grqq{} (\textit{coverage analysis}) in der medizinischen Informatik prim\"ar f\"ur die terminologische Zuordnung zwischen Klassifikationssystemen verwendet wird (z.\,B.\ ICF vs.\ SNOMED~CT f\"ur die Behinderungsbewertung). Das vorliegende Modell rekonzeptualisiert die Abdeckungsanalyse als klinische Qualit\"atssicherungsmetrik: Statt die \"Uberlappung zwischen Terminologien zu messen, quantifiziert sie den Anteil des Symptomprofils eines Patienten, der durch aktuelle Diagnosen erkl\"art wird.

\begin{figure}[ht]
\centering
\begin{tikzpicture}[
    font=\small
]
% Farben definieren
\definecolor{covgreen}{RGB}{76,153,0}
\definecolor{covyellow}{RGB}{204,163,0}
\definecolor{covred}{RGB}{204,51,51}

% Spaltenpositionen
\def\colA{3.2}   % Depression
\def\colB{5.2}   % PTBS
\def\colC{7.2}   % Hypothyreose
\def\colD{9.4}   % Abdeckung c(s_i)

% Zeilenpositionen (oben nach unten)
\def\rowH{6.0}   % Kopfzeile
\def\rowA{5.2}   % Niedergeschlagenheit
\def\rowB{4.4}   % Fatigue
\def\rowC{3.6}   % Insomnie
\def\rowD{2.8}   % Intrusionen
\def\rowE{2.0}   % Hyperarousal
\def\rowF{1.2}   % Konzentration
\def\rowT{0.2}   % Gesamt

% Kopfzeile
\node[font=\small\bfseries, anchor=west] at (0,\rowH) {Symptom};
\node[font=\small\bfseries, rotate=35, anchor=south west] at (\colA-0.2,\rowH+0.1) {Depression};
\node[font=\small\bfseries, rotate=35, anchor=south west] at (\colB-0.2,\rowH+0.1) {PTBS};
\node[font=\small\bfseries, rotate=35, anchor=south west] at (\colC-0.2,\rowH+0.1) {Hypothyreose};
\node[font=\small\bfseries, rotate=35, anchor=south west] at (\colD-0.2,\rowH+0.1) {$c(s_i)$};

% Horizontale Linien
\draw[thick] (-0.1,\rowH-0.25) -- (\colD+0.7,\rowH-0.25);
\draw[thick] (-0.1,\rowT-0.3) -- (\colD+0.7,\rowT-0.3);
\draw[thick] (-0.1,\rowT+0.45) -- (\colD+0.7,\rowT+0.45);

% Zeile A: Niedergeschlagenheit
\node[anchor=west] at (0,\rowA) {Niedergeschl.};
\node at (\colA,\rowA) {0.90};
\node at (\colB,\rowA) {0.30};
\node at (\colC,\rowA) {0.20};
\fill[covgreen, opacity=0.25, rounded corners=2pt] (\colD-0.45,\rowA-0.25) rectangle (\colD+0.45,\rowA+0.25);
\node[font=\small\bfseries] at (\colD,\rowA) {0.90};

% Zeile B: Fatigue
\node[anchor=west] at (0,\rowB) {Fatigue};
\node at (\colA,\rowB) {0.70};
\node at (\colB,\rowB) {0.20};
\node at (\colC,\rowB) {0.80};
\fill[covyellow, opacity=0.25, rounded corners=2pt] (\colD-0.45,\rowB-0.25) rectangle (\colD+0.45,\rowB+0.25);
\node[font=\small\bfseries] at (\colD,\rowB) {0.80};

% Zeile C: Insomnie
\node[anchor=west] at (0,\rowC) {Insomnie};
\node at (\colA,\rowC) {0.80};
\node at (\colB,\rowC) {0.50};
\node at (\colC,\rowC) {0.10};
\fill[covyellow, opacity=0.25, rounded corners=2pt] (\colD-0.45,\rowC-0.25) rectangle (\colD+0.45,\rowC+0.25);
\node[font=\small\bfseries] at (\colD,\rowC) {0.80};

% Zeile D: Intrusionen
\node[anchor=west] at (0,\rowD) {Intrusionen};
\node at (\colA,\rowD) {0.10};
\node at (\colB,\rowD) {0.95};
\node at (\colC,\rowD) {0.00};
\fill[covgreen, opacity=0.25, rounded corners=2pt] (\colD-0.45,\rowD-0.25) rectangle (\colD+0.45,\rowD+0.25);
\node[font=\small\bfseries] at (\colD,\rowD) {0.95};

% Zeile E: Hyperarousal
\node[anchor=west] at (0,\rowE) {Hyperarousal};
\node at (\colA,\rowE) {0.15};
\node at (\colB,\rowE) {0.90};
\node at (\colC,\rowE) {0.05};
\fill[covgreen, opacity=0.25, rounded corners=2pt] (\colD-0.45,\rowE-0.25) rectangle (\colD+0.45,\rowE+0.25);
\node[font=\small\bfseries] at (\colD,\rowE) {0.90};

% Zeile F: Konzentration
\node[anchor=west] at (0,\rowF) {Konzentration};
\node at (\colA,\rowF) {0.65};
\node at (\colB,\rowF) {0.40};
\node at (\colC,\rowF) {0.55};
\fill[covyellow, opacity=0.25, rounded corners=2pt] (\colD-0.45,\rowF-0.25) rectangle (\colD+0.45,\rowF+0.25);
\node[font=\small\bfseries] at (\colD,\rowF) {0.65};

% Gesamtzeile
\node[anchor=west, font=\small\bfseries] at (0,\rowT) {$C_{\text{total}}$};
\fill[covgreen, opacity=0.25, rounded corners=2pt] (\colD-0.55,\rowT-0.25) rectangle (\colD+0.55,\rowT+0.25);
\node[font=\small\bfseries] at (\colD,\rowT) {83.3\%};

% Legende
\node[anchor=west, font=\scriptsize] at (0,-0.5) {%
    \tikz\fill[covgreen, opacity=0.4] (0,0) rectangle (0.3,0.25); \(\geq\)85\% (vollst.) \quad
    \tikz\fill[covyellow, opacity=0.4] (0,0) rectangle (0.3,0.25); 60--84\% (partiell) \quad
    \tikz\fill[covred, opacity=0.4] (0,0) rectangle (0.3,0.25); $<$60\% (unzureich.)
};

\end{tikzpicture}
\caption{Illustrative Abdeckungsanalyse-Matrix. Jede Zelle zeigt das Gewicht $w(s_i, d_j)$; die Abdeckungsspalte zeigt $c(s_i) = \max_j w(s_i, d_j)$. Der Gesamtabdeckungswert $C_{\text{total}} = 83{,}3\%$ deutet auf partielle Gesamtabdeckung hin und legt weitere diagnostische Untersuchung nahe.}
\label{fig:coverage}
\end{figure}

Die Abdeckungsanalyse wird symmetrisch auch in Achse~III implementiert (Subachse~IIIi): K\"orpersymptome, die durch keine aktuelle somatische Diagnose erkl\"art werden, werden systematisch identifiziert und f\"uhren zum medizinischen Untersuchungsplan (Subachse~IIIj).


% ===================================================================
% 7. DIMENSIONALE INTEGRATION
% ===================================================================
\section{Dimensionale Integration: HiTOP und RDoC als erg\"anzende Perspektiven}
\label{sec:dimensional}

W\"ahrend die Abdeckungsanalyse die Vollst\"andigkeit \textit{innerhalb} der kategorialen Diagnostik quantifiziert, adressieren dimensionale Rahmenmodelle eine komplement\"are Limitation: die k\"unstliche Grenzziehung zwischen diagnostischen Kategorien \citep{Rosenman2003}. Das 6-Achsen-Modell basiert prim\"ar auf kategorialer Diagnostik (DSM-5-TR/ICD-11), enth\"alt aber bereits dimensionale Elemente (PID-5, WHODAS~2.0). Zwei weitere dimensionale Rahmenmodelle verdienen Integration als erg\"anzende Schichten:

\subsection{HiTOP (Hierarchical Taxonomy of Psychopathology)}

Die Hierarchical Taxonomy of Psychopathology \citep{Kotov2017, Kotov2022, Krueger2018} organisiert psychische St\"orungen empirisch-hierarchisch in sechs Spektren: Internalizing (Depression, Angst, PTBS), Thought Disorder (Psychose, Manie), Disinhibited Externalizing (Substanzgebrauch, Impulskontrolle), Antagonistic Externalizing (Antisoziales Verhalten), Detachment (soziale Zur\"uckgezogenheit, Anhedonie) und Somatoform (somatische Symptome). Die Cross-Cutting-Ergebnisse des Systems werden direkt den HiTOP-Spektren zugeordnet und als Radardiagramm visualisiert. Die Zuordnung nutzt die maximale Auspr\"agung der zugeh\"origen Cross-Cutting-Dom\"anen:

\begin{itemize}[noitemsep]
\item Internalizing = max(Depression, Anxiety, Sleep)
\item Thought Disorder = max(Psychosis, Mania, Dissociation)
\item Disinhibited Externalizing = max(Substance)
\item Antagonistic Externalizing = max(Anger)
\item Detachment = max(Personality/Withdrawal)
\item Somatoform = max(Somatic)
\end{itemize}

\textbf{Anmerkung zur HiTOP-Zuordnung:} Die Zuordnung folgt dem aktuellen HiTOP-Konsensus \citep{Kotov2021}: Manie wird dem Thought-Disorder-Spektrum zugeordnet (nicht Disinhibited Externalizing), da manische Episoden in der hierarchischen Faktorstruktur prim\"ar mit dem Thought-Disorder-Faktor laden. Dissoziation wird ebenfalls Thought Disorder zugeordnet, konsistent mit der empirischen Evidenz f\"ur eine enge Assoziation mit psychotischen Erfahrungen. Das Somatoform-Spektrum verwendet ausschlie\ss{}lich die somatische Cross-Cutting-Dom\"ane; die Depression-assoziierte somatische Belastung wird \"uber das Internalizing-Spektrum erfasst. Diese Zuordnungen sind als Approximationen zu verstehen, da die Cross-Cutting-Dom\"anen nicht f\"ur die HiTOP-Zuordnung entwickelt wurden.

Diese automatische Berechnung erfordert \textit{keine zus\"atzliche Datenerhebung} --- die bereits im Cross-Cutting-Screening (Abschnitt~\ref{sec:screening}) erfassten Dom\"anenwerte werden direkt wiederverwendet. Es ist jedoch zu betonen, dass diese Zuordnung eine \textit{grobe Approximation} darstellt: Die Cross-Cutting-Dom\"anen wurden nicht f\"ur die HiTOP-Zuordnung entwickelt, und die Max-Funktion bildet die faktorenanalytisch ermittelten Spektren nur unvollst\"andig ab. Die Darstellung als Radardiagramm (visuell klar vom PID-5-Pers\"onlichkeitsprofil abgesetzt) kann die Erkennung transdiagnostischer Muster f\"ordern, die bei rein kategorialer Diagnostik unsichtbar bleiben, sollte aber als explorative Visualisierung, nicht als validierte Messung interpretiert werden.

Zusammengefasst bildet die beschriebene Zuordnung eine erste Approximation, die mit zunehmender empirischer Evidenz zu HiTOP-konformen Assessments verfeinert werden sollte \citep{Kotov2021}.

\subsection{RDoC (Research Domain Criteria)}

Die Research Domain Criteria des NIMH \citep{Insel2010} definieren sechs Dom\"anen (Negative Valence, Positive Valence, Cognitive Systems, Social Processes, Arousal/Regulatory, Sensorimotor) mit sieben Analyseebenen (Gene, Molek\"ule, Zellen, Schaltkreise, Physiologie, Verhalten, Selbstbericht). F\"ur ein klinisches System ist RDoC prim\"ar als \textit{Forschungsannotation} relevant: Wenn Biomarker-Daten verf\"ugbar sind (z.\,B.\ Cortisol-Profile, Neuroimaging), k\"onnen diese den RDoC-Dom\"anen zugeordnet werden. Das APA Future DSM Strategic Committee exploriert explizit die Integration von Biomarkern --- das System ist darauf vorbereitet.

Ein separates Konzeptpapier beschreibt die detaillierte Architektur der dimensionalen Integration.


% ===================================================================
% 8. TECHNISCHE ARCHITEKTUR
% ===================================================================
\section{Technische Architektur der Python-Implementierung}
\label{sec:technik}

\subsection{Hierarchische Zustandsmaschinen als Entscheidungsmotor}

Nach Evaluation von vier Architekturans\"atzen --- AnyTree, Zustandsmaschinen, Rule Engines und Behavior Trees --- erweisen sich \textbf{Hierarchische Zustandsmaschinen (HSMs)} mittels der Python-Bibliothek \texttt{transitions} als optimale L\"osung f\"ur psychiatrische Diagnostik-Workflows.

Die \texttt{transitions}-Bibliothek mit ihrer \texttt{HierarchicalMachine}-Erweiterung bietet: konditionelle Transitionen via Guards (direkte Abbildung von \glqq wenn Symptom~X $\to$ betrete Modul~Y\grqq{}), verschachtelte Zust\"ande (der 6-Stufen-Prozess mit st\"orungsspezifischen Submodulen), History States (R\"ucknavigation) und serialisierbaren Zustand (Speichern/Fortsetzen).

Die Architektur bildet sich wie folgt ab:

\begin{verbatim}
Top-level: [Intake -> Step1_Malingering -> Step2_Substance ->
            Step3_Medical -> Step4_CrossCutting ->
            Step5_DisorderModules -> Step6_Functioning -> Summary]

Step5_DisorderModules (verschachtelt, 11 Module):
  +-- MoodDisorders -> {MDD_Criteria, Bipolar_Screening, Dysthymia,
                        PMDD, Prolonged_Grief}
  +-- AnxietyDisorders -> {GAD, Panic, Social_Anxiety, Phobias,
                           Separation_Anxiety, Selective_Mutism}
  +-- TraumaDisorders -> {PTSD_DSM5, PTSD_ICD11, CPTSD_DSO,
                          Acute_Stress, Adjustment}
  +-- PsychoticDisorders -> {Schizophrenia_1mo, Schizophrenia_6mo,
                             Schizoaffective, Brief_Psychotic}
  +-- PersonalityDisorders -> {LPFS, PID5_Traits, ICD11_Severity}
  +-- OCDSpectrum -> {OCD_Criteria, BDD, Hoarding, Trichotillomania,
                      Excoriation}
  +-- DissociativeDisorders -> {DID, Depersonalization,
                                Dissociative_Amnesia}
  +-- EatingDisorders -> {AN, BN, BED, ARFID}
  +-- NeurodevelopmentalDisorders -> {ADHD, ASD_Assessment}
  +-- SubstanceUseDisorders -> {Alcohol_Use, Drug_Use,
                                Behavioral_Addictions}
  +-- SomaticDisorders -> {Somatic_Symptom, Illness_Anxiety,
                           Conversion, Factitious}
\end{verbatim}

Die Erweiterung von 5 auf \textbf{11~St\"orungsmodule} stellt sicher, dass jedes Screening-Instrument der Matrix (Tabelle~\ref{tab:instrumente}) in einen diagnostischen Workflow m\"undet. Ein auff\"alliger DES-II-Score triggert das Modul DissociativeDisorders; ein erh\"ohter SCOFF-Score das Modul EatingDisorders. Ohne diese Vollst\"andigkeit best\"unde eine systemische L\"ucke zwischen Screening und Diagnostik.

AnyTree wird weiterhin f\"ur \textbf{Visualisierung und Serialisierung} der Baumstruktur genutzt (Graphviz-Export, JSON-Repr\"asentation), w\"ahrend \texttt{transitions} die Laufzeitlogik verwaltet.

\subsection{Empfohlener Technologie-Stack}

\begin{table}[H]
\centering
\caption{Empfohlener Technologie-Stack}
\label{tab:techstack}
\begin{tabularx}{\textwidth}{lXl}
\toprule
\textbf{Komponente} & \textbf{Empfehlung} & \textbf{Begr\"undung} \\
\midrule
Entscheidungsmotor & \texttt{transitions} (HSM) & Verschachtelte Zust\"ande, Guards \\
Baumvisualisierung & \texttt{anytree} & Graphviz-Export, JSON \\
UI (Prototyp) & Streamlit & Schnelles Prototyping, \texttt{st.navigation} \\
PDF-Berichte & WeasyPrint + Jinja2 & Natives UTF-8, HTML/CSS-Templates \\
Pers\"onlichkeits-Charts & Plotly & \texttt{px.line\_polar()} f\"ur PID-5-Radardiagramme \\
Datenpersistenz & SQLModel + SQLite & Pydantic + SQLAlchemy kombiniert \\
Datenvalidierung & Pydantic & Typsichere Diagnosekriterien-Modelle \\
ICD-11-Codes & WHO ICD-11 API & Strukturierte Diagnose-Code-Suche \\
Interoperabilit\"at & HL7 FHIR (R4) & Standardisierter Datenaustausch \\
\bottomrule
\end{tabularx}
\end{table}

\subsection{Anbindung der Klassifikationskataloge}
\label{sec:katalog-anbindung}

Ein Werkzeug zur diagnostischen Codierung muss eine Frage beantworten, bevor es irgendeine klinische beantworten kann: \emph{Woher stammt dieser Code?} Die Referenzimplementierung beantwortet sie, indem sie die Klassifikationen gar nicht erst selbst besitzt.

Frühere Versionen lieferten einen handkuratierten Seed-Katalog psychiatrischer Codes im Repository mit. Für die Demonstration eines Achsenmodells genügt das, aber es ist die typische Fehlerform klinischer Codiersoftware: Eine private Kopie einer Klassifikation läuft still von der gepflegten Fassung weg, und die Abweichung bleibt genau dort unsichtbar, wo sie zählt --- im Moment des Codierens. Die Implementierung pflegt deshalb keine Kopie mehr. Sie \textbf{bindet externe Katalogmodule} hinter einer einheitlichen Provider-Schnittstelle an (Tabelle~\ref{tab:katalog-anbindung}); der mitgelieferte Seed-Katalog wird zum ausdrücklichen Offline-Rückfallpfad zurückgestuft.

\begin{table}[H]
\centering
\caption{Angebundene externe Klassifikationskataloge}
\label{tab:katalog-anbindung}
\begin{tabularx}{\textwidth}{lXll}
\toprule
\textbf{Katalog} & \textbf{Angebunden über} & \textbf{Datenlizenz} & \textbf{Betriebsart} \\
\midrule
ICD-10 (WHO 2019) & \texttt{simple-icd-10} & Public Domain (CC0) & Offline \\
ICD-10-CM (USA) & \texttt{simple-icd-10\allowbreak-cm} & Public Domain (CDC/NCHS) & Offline \\
ICD-11 (WHO) & \texttt{simple-icd-11} $\rightarrow$ WHO ICD-API & CC BY-ND 3.0 IGO & Online, zugangsgebunden \\
\bottomrule
\end{tabularx}
\end{table}

Zwei Eigenschaften der Anbindung sind bewusste Entwurfsentscheidungen und keine Implementierungsdetails.

\textbf{Die Herkunft ist explizit.} Jeder aufgelöste Code führt Quelle, Lizenz, gebundenes Modul und Stand mit sich, und die Oberfläche zeigt an, welcher Katalog tatsächlich geantwortet hat. Ein Codiervorschlag, dessen Herkunft nicht benennbar ist, ist kein brauchbarer Codiervorschlag; ein unbeschrifteter Code in einer Krankenakte ist eher ein Haftungsrisiko als eine Hilfe.

\textbf{Die Verifikation ist dreiwertig.} Eine Abfrage liefert \texttt{verified}, \texttt{not\_found} oder \texttt{unavailable}. Der dritte Zustand existiert, weil der Unterschied zwischen \emph{\glqq der Katalog sagt, diesen Code gibt es nicht\grqq{}} und \emph{\glqq wir konnten den Katalog nicht fragen\grqq{}} klinisch tragend ist. Würde man \texttt{unavailable} auf \texttt{not\_found} abbilden, erschiene ein nicht überprüfbarer Code als widerlegt; die Abbildung auf \texttt{verified} wäre noch schlimmer. Das System tut weder das eine noch das andere.

Am deutlichsten wird das bei ICD-11. Die WHO-ICD-API ist die maßgebliche Quelle und verlangt OAuth2-Zugangsdaten (kostenlose Registrierung unter \texttt{icd.who.int/icdapi}). Fehlen Zugangsdaten oder Netzverbindung, meldet sich der ICD-11-Provider als nicht verfügbar, und die Anwendung fällt auf den Seed-Katalog zurück --- sie erfindet keine ICD-11-Codes, um die Lücke zu füllen. Die Lizenz CC BY-ND 3.0 IGO auf den WHO-ICD-11-Inhalten ist ein weiterer Grund, abzufragen statt mitzuliefern: Die No-Derivatives-Klausel verbietet die Weitergabe des Katalogs in veränderter Form, und genau das wäre eine lokal transformierte Kopie.

Die im Validierungs-Gate-Ledger (Abschnitt~\ref{sec:validierung}) gezogene Grenze gilt für diese Anbindungen unverändert. Eine korrekte Katalogabfrage belegt \emph{Codier- und Dokumentationsgenauigkeit}. Sie belegt keine diagnostische Validität, und keine noch so große Zahl korrekt aufgelöster Codes darf als klinische Validierung des Achsenmodells ausgegeben werden.

\subsection{Open-Source-Referenzimplementierung}

Der vollst\"andige Quellcode des Prototyps (V9, ca.\ 1.850~Zeilen Python/Streamlit), die bilinguale \"Ubersetzungsdatei, der Entwicklungsfahrplan und das vorliegende Manuskript sind \"offentlich verf\"ugbar unter:

\begin{center}
\href{https://github.com/um-bruch/multiaxial-diagnostic-system}{github.com/um-bruch/multiaxial-diagnostic-system}
\end{center}

\noindent Die Bereitstellung als Open-Source-Repository dient der wissenschaftlichen Transparenz und Reproduzierbarkeit. Installationsanweisungen finden sich in der \texttt{README.md} des Repositories.

\textbf{Lizenzierung und Archivierung:} Das Repository steht unter der MIT-Lizenz, die sowohl akademische als auch kommerzielle Nutzung erlaubt, jedoch keinerlei Gew\"ahrleistung \"ubernimmt (vgl.\ Haftungsfragen, Abschnitt~\ref{sec:ethik}). F\"ur die langfristige Zitierbarkeit und Versionierung ist das Repository bereits als Software-Record auf Zenodo archiviert und mit dem Beitrag verkn\"upft (Concept-DOI: \href{https://doi.org/10.5281/zenodo.18736725}{10.5281/zenodo.18736725}; neueste \"offentliche Version zum Zeitpunkt dieser Manuskriptaktualisierung: \href{https://doi.org/10.5281/zenodo.19073268}{10.5281/zenodo.19073268}). Dadurch bleibt eine konkrete Codeversion dauerhaft referenzierbar, unabh\"angig von \"Anderungen am GitHub-Repository.


% ===================================================================
% 9. FUNKTIONALE ASSESSMENT-INTEGRATION
% ===================================================================
\section{Funktionale Beurteilung: Br\"ucke zwischen Diagnose und Behinderung}
\label{sec:funktional}

Die Integration funktionaler Beurteilung verbindet Diagnose und Teilhabe. Das Modell implementiert drei Ebenen: das krankheitsspezifische Assessment (ICF Core Sets), die allgemeine Behinderungsmessung (WHODAS~2.0) und die rechtlich-administrative Bewertung (GdB).

ICD-11 selbst hat sich in Richtung ICF-Integration bewegt, indem \glqq Functioning Properties\grqq{} eingef\"uhrt wurden, die mit 103 rehabilitationsrelevanten Gesundheitszust\"anden verkn\"upft sind. F\"ur Schizophrenie und psychotische St\"orungen enth\"alt ICD-11 \textbf{dimensionale Qualifikatoren}, die mit rehabilitationsbasierter psychiatrischer Versorgung konsistent sind.


% ===================================================================
% 14a. KLINISCHE FALLVIGNETTE
% ===================================================================
\section{Beispielhafte Fallvignette}
\label{sec:fallvignette}

\textit{Hinweis: Die folgende Fallvignette ist fiktiv und dient ausschlie\ss{}lich der Illustration der Anwendbarkeit des 6-Achsen-Modells. Jede \"Ahnlichkeit mit realen Personen ist zuf\"allig.}

\vspace{0.5em}
\noindent \textbf{Patient M., 34~Jahre, m\"annlich, Erstvorstellung in der psychiatrischen Ambulanz.}

\begin{description}[style=nextline, leftmargin=1.5cm, font=\bfseries\small]
\item[Achse~I (Psychische Profile)]
Ia (akut): Mittelgradige depressive Episode (F32.1), Konfidenz 85\%, PRO: PHQ-9 = 16, Antriebslosigkeit, Schlafst\"orungen seit 4~Monaten; CONTRA: keine Suizidalit\"at, erhaltene Arbeitsf\"ahigkeit. \\
Ih (Verdacht): Generalisierte Angstst\"orung (F41.1), Konfidenz 55\%, PRO: GAD-7 = 12, chronische Sorgen; CONTRA: m\"oglicherweise sekund\"ar zur Depression. \\
Ii (Abdeckung): Gesamtabdeckung $\sim$78\%. Nicht abgedeckt: chronische R\"uckenschmerzen (somatisch? funktionell?), Konzentrationsprobleme (depressionsassoziiert oder eigenst\"andig?). \\
Ij (Plan): Wichtig: Neuropsychologische Testung zur Kl\"arung der Konzentrationsst\"orung; Wichtig: Somatische Abkl\"arung der R\"uckenschmerzen (funktionell vs.\ strukturell); Verlauf: PHQ-9-Monitoring alle 4~Wochen, C-SSRS-Screening bei Verschlechterung.
\item[Achse~II (Biographie)]
PID-5-BF+M: Erh\"ohte Negative Affektivit\"at (T=68), grenzwertige Distanziertheit (T=60). \\
Pr\"agende Erfahrung: Fr\"uhe Scheidung der Eltern (Alter 7), intermittierendes Mobbing in der Schule (12--15~J.). \\
Grundkonflikt: Autonomie-Abh\"angigkeits-Konflikt (OPD).
\item[Achse~III (Med. Synopse)]
IIIa: Lumbales Schmerzsyndrom (M54.5). IIIc: Hypothyreose (E03.9), beitragend zur Fatigue --- Kausalanalyse (IIIk): TSH = 6.2~mU/L, partielle Erkl\"arung der Antriebsminderung. \\
IIIm: L-Thyroxin 75\,$\mu$g, Wirkung 6/10; Ibuprofen 400\,mg b.B., Wirkung 5/10; keine relevanten Interaktionen.
\item[Achse~IV (Umwelt \& Funktion)]
WHODAS~2.0 Gesamtscore: 28/100 (moderate Beeintr\"achtigung). Mini-ICF-APP: Einschr\"ankungen in Regelm\"a\ss{}igkeit (3/4), Kontaktf\"ahigkeit (2/4). GdB: nicht beantragt. \\
Stressoren: Drohender Arbeitsplatzverlust, Partnerschaftskrise. \\
Bezugspersonen: Partnerin (ambivalent st\"utzend), Hausarzt Dr.~X, kein Therapeut bisher.
\item[Achse~V (Bedingungsmodell)]
P1 (Pr\"adisposition): Fr\"uhe Verlusterfahrung, erh\"ohte Negative Affektivit\"at. \\
P2 (Ausl\"oser): Arbeitsplatzbedrohung seit 4~Monaten, zeitlich konkordant mit Symptombeginn. \\
P3 (Aufrechterhaltung): Sozialer R\"uckzug, unbehandelte Hypothyreose, Schmerzchronifizierung. \\
P4 (Protektiv): Durchschnittliche Intelligenz, stabile Wohnsituation, Problembewusstsein.
\item[Achse~VI (Belegsammlung)]
Evidenz-Matrix: PHQ-9 (21.02.2026), GAD-7 (21.02.2026), Labor (TSH, CRP --- 15.02.2026). \\
CAVE: Hypothyreose beitragende Ursache der Fatigue --- Schilddr\"usenparameter vor Antidepressiva-Einstellung kontrollieren! \\
Symptomverlauf: Antriebslosigkeit seit November 2025, progredient; R\"uckenschmerzen seit 2024, fluktuierend.
\end{description}

\noindent Diese Vignette illustriert, wie das 6-Achsen-Modell eine integrierte diagnostische Perspektive erm\"oglicht: Die Abdeckungsanalyse (78\%) --- berechnet als $C_{\text{total}} = \frac{1}{n} \sum c(s_i)$ \"uber die Symptome Niedergeschlagenheit, Antriebslosigkeit, Schlafst\"orungen, Sorgen, R\"uckenschmerzen und Konzentrationsst\"orungen, wobei R\"uckenschmerzen und Konzentrationsst\"orungen nur partiell durch bestehende Diagnosen abgedeckt sind (vgl.\ Abschnitt~\ref{sec:abdeckung}) --- identifiziert ungekl\"arte Symptome. Der CAVE-Hinweis verhindert eine vorschnelle Antidepressiva-Einstellung ohne Schilddr\"usenkontrolle, und das Bedingungsmodell verbindet biographische Vulnerabilit\"at mit aktuellem Ausl\"oser.


% ===================================================================
% 10. ETHIK UND DATENSCHUTZ
% ===================================================================
\section{Ethik, Datenschutz und algorithmischer Bias}
\label{sec:ethik}

Die computergest\"utzte psychiatrische Diagnostik wirft spezifische ethische Fragen auf, die bei der Entwicklung und Implementierung systematisch adressiert werden m\"ussen.

\subsection{Algorithmischer Bias und Validierungspopulationen}

Screening-Instrumente sind an spezifischen Populationen validiert, und die \"Ubertragbarkeit ihrer Cutoff-Werte auf andere Populationen ist limitiert. Der PHQ-9 zeigt kulturell variable Cutoff-Werte: W\"ahrend f\"ur westliche Populationen ein Cutoff von $\geq 10$ validiert ist, liegen in asiatischen L\"andern h\"aufig niedrigere optimale Schwellenwerte vor. Der AQ-10 hat eine bekannte Geschlechterverzerrung bei Frauen mit Autismus-Spektrum-St\"orungen (\"Uberdiagnose bei M\"annern, Unterdiagnose bei Frauen aufgrund unterschiedlicher Ph\"anotypen). F\"ur das AUDIT (Alkoholscreening) werden in der Literatur niedrigere geschlechtsspezifische Cutoffs f\"ur Frauen diskutiert, da die Standardschwelle von $\geq 8$ m\"oglicherweise Alkoholprobleme bei Frauen untersch\"atzt \citep{Saunders1993}.

\textbf{Sozio\"okonomischer Bias:} Instrumente wie das WHODAS~2.0 erfassen Funktionseinschr\"ankungen, die bei Personen mit niedrigem sozio\"okonomischem Status teilweise durch fehlende Ressourcen (z.\,B.\ Transport, Kinderbetreuung) bedingt sind, nicht durch die psychische St\"orung selbst. Das System muss diese strukturellen Faktoren in Achse~IV (psychosoziale Umst\"ande) separat dokumentieren, um Fehlattributionen zu vermeiden.

\textbf{Migrationshintergrund und Sprachbarrieren:} Selbstberichts-Instrumente setzen ausreichende Sprachkompetenz voraus. Bei Personen mit Migrationshintergrund k\"onnen sprachliche Missverst\"andnisse zu falsch-positiven oder falsch-negativen Ergebnissen f\"uhren. Das System sollte bei nicht-muttersprachlichen Patienten einen expliziten Hinweis auf diese Limitation ausgeben.

\textbf{Mitigation:} Das System muss diese Bias-Risiken transparent dokumentieren und --- wo verf\"ugbar --- populationsspezifische Normen anbieten. Die Integration des Cultural Formulation Interview (Abschnitt~\ref{sec:achse4}) adressiert kulturellen Bias auf der Systemebene. F\"ur jedes Screening-Instrument sollte die Validierungspopulation dokumentiert sein, sodass Kliniker die \"Ubertragbarkeit auf den konkreten Fall beurteilen k\"onnen.

\subsection{Datenschutz und DSGVO}

Psychiatrische Diagnosedaten geh\"oren zu den sensibelsten Gesundheitsdaten. Die Implementierung muss folgende Anforderungen erf\"ullen: Verschl\"usselung aller gespeicherten und \"ubertragenen Daten (AES-256, TLS~1.3); optionale Zugriffskontrolle nach Professionszuordnung (Abschnitt~\ref{sec:multiprofessionell}); Audit-Trail f\"ur alle diagnostischen Entscheidungen; Recht auf L\"oschung und Datenportabilit\"at; Datensparsamkeit --- nur diagnostisch notwendige Informationen werden erfasst. Besondere Vorsicht gilt f\"ur die Belegsammlung (Achse~VI), die pers\"onliche Dokumente referenziert.

\subsection{Informierte Einwilligung und Patientenautonomie}

Die Nutzung eines computergest\"utzten Diagnosesystems wirft Fragen der informierten Einwilligung auf. Patienten haben das Recht zu wissen, dass ihre diagnostische Beurteilung teilweise durch algorithmische Verfahren unterst\"utzt wird. Das System sollte daher folgende Transparenzanforderungen erf\"ullen:

\begin{itemize}[noitemsep]
\item \textbf{Aufkl\"arung \"uber Systemnutzung:} Patienten werden dar\"uber informiert, dass Screening-Instrumente und diagnostische Entscheidungsb\"aume als Unterst\"utzungswerkzeuge eingesetzt werden.
\item \textbf{Opt-out-M\"oglichkeit:} Patienten haben das Recht, eine rein menschlich durchgef\"uhrte Diagnostik zu verlangen, falls sie der computergest\"utzten Methode nicht vertrauen.
\item \textbf{Erkl\"arbarkeit von Diagnosevorschl\"agen:} Das System muss nachvollziehbar machen, welche Symptome und Kriterien zu einem Diagnosevorschlag gef\"uhrt haben. Die PRO/CONTRA-Evidenzbewertung (Abschnitt~\ref{sec:achse1}) dient genau diesem Zweck.
\item \textbf{Zugang zu eigenen Daten:} Patienten haben das Recht auf Einsicht in ihre diagnostische Akte (Achse~I--VI) und auf Erkl\"arung der dokumentierten Befunde.
\end{itemize}

Diese Anforderungen sind nicht nur ethisch geboten, sondern auch rechtlich durch die DSGVO (Art.~13--15) verankert.

\subsection{Haftungsfragen und medizinisch-rechtliche Verantwortung}

Die Nutzung eines computergest\"utzten Diagnosesystems verschiebt nicht die Haftung vom Kliniker auf das System. Folgende rechtliche Prinzipien gelten:

\begin{description}[style=nextline, leftmargin=1.5cm, font=\bfseries]
\item[Behandler-Haftung bleibt bestehen] Der behandelnde Arzt oder Psychotherapeut bleibt vollumf\"anglich haftbar f\"ur diagnostische Fehlentscheidungen, auch wenn diese durch ein algorithmisches System unterst\"utzt wurden. Das System ist ein Werkzeug, kein Entscheidungstr\"ager.
\item[Sorgfaltspflicht bei Systemnutzung] Kliniker m\"ussen das System sachgem\"a\ss{} nutzen und dessen Grenzen kennen. Eine blinde \"Ubernahme algorithmischer Diagnosevorschl\"age ohne klinische Pr\"ufung erf\"ullt nicht die \"arztliche Sorgfaltspflicht.
\item[Dokumentationspflicht] Die Nutzung des Systems und die klinische \"Uberpr\"ufung der Systemvorschl\"age sollten dokumentiert werden (Achse~VI, Kontaktprotokoll). Dies dient der Nachvollziehbarkeit im Haftungsfall.
\item[Produkthaftung des Entwicklers] Der Systementwickler haftet f\"ur nachweisliche Softwarefehler (z.\,B.\ falsche Implementierung von DSM-5-TR-Kriterien, Rechenfehler in der Abdeckungsanalyse). Die Open-Source-Ver\"offentlichung (Abschnitt~\ref{sec:technik}) erm\"oglicht externe Code-Reviews und erh\"oht die Transparenz.
\item[Medizinprodukt-Einstufung bedarf juristischer Kl\"arung] Nach EU-Medizinprodukteverordnung (MDR 2017/745) sind reine Dokumentationssysteme ohne therapeutische oder pr\"adiktive Funktion nicht zwingend als Medizinprodukt einzustufen. Das vorliegende System geht jedoch \"uber reine Dokumentation hinaus: Die Gatekeeper-Logik, die algorithmischen Entscheidungsb\"aume und die Diagnosevorschl\"age k\"onnten als Entscheidungsunterst\"utzungssoftware im Sinne der MDR klassifiziert werden. Eine juristische Pr\"ufung dieser Einstufung ist daher \textit{zwingend erforderlich}, bevor das System in der klinischen Praxis eingesetzt wird.
\end{description}

Die klare Trennung zwischen \textit{diagnostischer Unterst\"utzung} (zul\"assig) und \textit{autonomer Diagnosestellung} (nicht implementiert) ist essentiell f\"ur die rechtliche Absicherung.

\subsection{Klinische Verantwortung und epistemische Grenzen}

Das System ist als \textit{Unterst\"utzungswerkzeug} konzipiert, nicht als diagnostischer Automat. Alle algorithmischen Diagnosevorschl\"age erfordern klinische Best\"atigung. Die epistemische Grenze der Automatisierbarkeit (vgl.\ Stufe~6 der Gatekeeper-Logik, Abschnitt~\ref{sec:gatekeeper}) wird im System explizit kommuniziert. Die finale diagnostische Verantwortung liegt stets beim behandelnden Kliniker.

\textbf{Kritische F\"alle, die algorithmische Systeme nicht zuverl\"assig bewerten k\"onnen:}
\begin{itemize}[noitemsep]
\item \textbf{Suizidalit\"at:} Selbstberichts-Instrumente (C-SSRS) sind hilfreich, aber jede Best\"atigung suizidaler Gedanken erfordert ein ausf\"uhrliches klinisches Gespr\"ach. Algorithmische Risikobewertungen k\"onnen die klinische Einsch\"atzung nicht ersetzen.
\item \textbf{Psychotische Symptome:} Selbstberichts-Instrumente (PQ-16) sind Screening-Werkzeuge, keine diagnostischen Instrumente. Eine Best\"atigung psychotischer Symptome erfordert klinische Exploration und oft Fremdanamnese.
\item \textbf{Kulturgebundene Syndrome:} Das DSM-5-TR enth\"alt Glossare kulturgebundener Syndrome (z.\,B.\ Ataque de nervios, Taijin kyofusho), die algorithmisch schwer zu erfassen sind. Hier ist kulturelle Kompetenz des Klinikers unersetzlich.
\end{itemize}

Diese Grenzen machen das System nicht unbrauchbar, sondern unterstreichen die Notwendigkeit einer \textit{hybriden} Diagnostik, in der algorithmische Unterst\"utzung und klinisches Urteil komplement\"ar wirken.


% ===================================================================
% 11. VALIDIERUNGSSTRATEGIE
% ===================================================================
\section{Validierungsstrategie}
\label{sec:validierung}

\textbf{Aktueller Status:} Der vorliegende Artikel pr\"asentiert die konzeptionelle Architektur und technische Implementierung des 6-Achsen-Modells. Die nachfolgend beschriebene Validierungsstrategie stellt einen \textit{Implementierungsplan} dar --- die empirische Validierung ist zum gegenw\"artigen Zeitpunkt noch nicht durchgef\"uhrt. Dieser Abschnitt beschreibt daher die geplanten Validierungsschritte, nicht bereits vorliegende Ergebnisse.

Die Abschaffung des DSM-IV-Systems wurde unter anderem durch mangelnde Interrater-Reliabilit\"at begr\"undet. Das vorliegende System muss daher von Anfang an eine rigorose Validierungsstrategie verfolgen:

\begin{table}[htbp]
\centering
\caption{Validierungs-Gate-Ledger f\"ur die geplante Post-Design-Evaluation. Die Zeilen definieren das minimale Datenobjekt, das vorliegen muss, bevor ein Designclaim zu einem Prozess-, Nutzen- oder klinischen Validit\"atsclaim angehoben werden darf.}
\scriptsize
\renewcommand{\arraystretch}{1.15}
\begin{tabularx}{\textwidth}{@{}L{2.7cm}Y Y Y@{}}
\toprule
Gate & Minimales Datenobjekt & Pass- / Hold-Regel & Verhinderter \"Uberclaim \\
\midrule
Diagnostisches Herleitungs-Ledger & Extrahierte Befunde, Differenzialhypothesen, Regel/Backing, Ausschlusskriterien, Unsicherheit und Gutachter-Notizen & Pass nur, wenn Prim\"ardiagnosen und verworfene Alternativen von Patientendaten zu klinischer Regel und Expertenbegutachtung r\"uckverfolgbar sind & Top-1-Diagnose\"ubereinstimmung oder LLM-Vorschlag wird als nachgewiesene Entscheidungsqualit\"at gelesen \\
Konsultations-Evidenzerhebungs-Gate & Frageprotokoll zu fehlenden DSM-/ICD-Kriterien, Dauer, Funktion, Fremdanamnese und somatischen Ausschl\"ussen & Hold, wenn ein fehlendes Kriterium aus Schweigen erschlossen statt aktiv erfragt oder als nicht verf\"ugbar markiert wird & Statische Vignettenkorrektheit wird als Konsultationskompetenz missverstanden \\
Synthetik-Fidelity-/Fairness-Gate & Vignettenmanifest mit Demografie, Basisratenannahmen, Schwellenverteilungen und Subgruppenfehlern & Pass nur, wenn synthetische F\"alle klinisch plausible Varianz erhalten und seltene oder Minderheitenprofile nicht kollabieren & Plausible Einzelf\"alle werden als populationsvalide Evaluationsdaten behandelt \\
KI-Agenten-Governance-Gate & Rollendefinition, Human-in-the-loop-Punkt, Eskalationsregel, Pr\"ufprotokoll, Datenschutzgrenze und \"Uberwachungsplan & Hold, wenn das System autonom diagnostisch handeln kann oder Verantwortung nicht zuordenbar ist & Ein Assistenzprototyp wird als klinisch validierte autonome KI dargestellt \\
ICD-Codierungsgrenze & Zuordnung der klinischen Entscheidung zu administrativem ICD-Code mit Codierrolle und Umgang mit Dissens & Pass nur, wenn die Codierung explizit nachgelagert bleibt und die klinische Diagnose nicht ver\"andert & Codiergenauigkeit wird als diagnostische Validit\"at gewertet \\
Achse-VI-Biomarker-Warteschlange & Signalprovenienz, Validierungspopulation, Informationsabfluss-Pr\"ufung, Subgruppenanalyse und prospektiver Best\"atigungsstatus & Hold, solange Wearable-, Stimm-, EMA- oder Chronosignale nicht repliziert sind und nur Kandidatensignale bleiben & Explorative digitale Signale werden zu diagnostischen Biomarkern erhoben \\
\bottomrule
\end{tabularx}
\label{tab:validierungs-gate-ledger}
\normalsize
\end{table}

Tabelle~\ref{tab:validierungs-gate-ledger} ist damit eine Designanforderungstabelle, keine Ergebnistabelle. Sie h\"alt die geplante Validierungssequenz an der zentralen Grenze des Modells fest: Strukturierte Entscheidungsunterst\"utzung kann vor klinischer Wirksamkeit evaluiert werden, darf aber ohne entsprechende Gate-Daten nicht als validierte klinische Diagnostik dargestellt werden.

\begin{description}[style=nextline, leftmargin=1.5cm, font=\bfseries]
\item[Phase~1: Expertenbegutachtung (N=10--15)] Strukturierte Begutachtung der Achsenarchitektur durch Psychiater, klinische Psychologen und Sozialarbeiter. Bewertung der Vollst\"andigkeit, klinischen Plausibilit\"at und Praktikabilit\"at. Ziel: Identifikation fehlender Subachsen oder \"uberfl\"ussiger Komponenten.

\item[Phase~2: Inter-Rater-Reliabilit\"atsstudie mit Fallvignetten (N=10 Rater, 20 Vignetten)] \textit{Dies ist die kritische Pilotstudie zur initialen Validierung des Systems:}

\textbf{Rekrutierung:} $N = 10$ unabh\"angige Rater (5 Psychiater, 5 approbierte Psychologische Psychotherapeuten; alle mit $\geq 3$ Jahren klinischer Erfahrung).

\textbf{Stimulusmaterial:} 20 standardisierte Fallvignetten (10 einfache F\"alle mit Einzeldiagnose, 10 komplexe F\"alle mit Komorbidit\"at), erstellt nach Vorbild der DSM-5-TR Clinical Cases \citep{BarnhillDSM5ClinicalCases}. Jede Vignette enth\"alt: Anamnese, aktuelle Beschwerden, Symptomverlauf, relevante medizinische Vorgeschichte, soziales Umfeld. Diagnosen vorab durch Expertenpanel (2 unabh\"angige Fach\"arzte + Konsensprozess) als Goldstandard etabliert.

\textbf{Prozedur:} Jeder Rater diagnostiziert alle 20 Vignetten (1) mit konventionellem Vorgehen (Freitext-Diagnose mit Begr\"undung), (2) mit dem 6-Achsen-System (vollst\"andige Achsendiagnostik). Reihenfolge der Methoden randomisiert; mindestens 1 Woche Washout zwischen Bedingungen pro Vignette.

\textbf{Prim\"are Outcome-Variablen:}
\begin{itemize}[leftmargin=*, noitemsep]
    \item \textit{Inter-Rater-Reliabilit\"at:} Fleiss' $\kappa$ f\"ur kategoriale Diagnosen (Achse~I Prim\"ardiagnose); Intraklassen-Korrelation (ICC) f\"ur dimensionale Ratings (PID-5-BF auf Achse~II, WHODAS~2.0 auf Achse~IV). Akzeptanzkriterium: $\kappa \geq 0{,}7$ (substanzielle \"Ubereinstimmung); ICC $\geq 0{,}75$.
    \item \textit{Konvergente Validit\"at:} \"Ubereinstimmung 6-Achsen-Diagnosen vs.\ Goldstandard. Sensitivit\"at und Spezifit\"at pro Diagnose.
    \item \textit{Inkrementelle Validit\"at der Abdeckungsanalyse:} Anteil der F\"alle, in denen Achse~Ii (unerkl\"arte Symptome) klinisch relevante Differenzialdiagnosen identifiziert, die im Freitext-Vorgehen \"ubersehen wurden. Hypothese: $\geq 20$\,\% der komplexen F\"alle profitieren von Abdeckungsanalyse.
    \item \textit{Bearbeitungszeit:} Durchschnittliche Bearbeitungszeit pro Vignette (konventionell vs.\ 6-Achsen).
\end{itemize}

\textbf{Sekund\"are Outcome-Variablen:}
\begin{itemize}[leftmargin=*, noitemsep]
    \item System Usability Scale (SUS) nach Abschluss aller Vignetten
    \item Subjektive Einsch\"atzung der N\"utzlichkeit einzelner Achsen (5-Punkt-Likert)
    \item Qualitative Interviews zu identifizierten Schwachstellen des Systems
\end{itemize}

\textbf{Power-Analyse:} Bei $N = 10$ Ratern und $k = 20$ Vignetten (Targets) kann ein ICC $= 0{,}75$ mit Power $1 - \beta = 0{,}8$ und $\alpha = 0{,}05$ gegen die Nullhypothese ICC $= 0{,}5$ getestet werden. Die pr\"azise KI-Breite h\"angt von der Anzahl der Targets ab, nicht von der Gesamtzahl der Urteile ($N \times k = 200$), da diese nicht unabh\"angig sind. Bei $k = 20$ Targets und $N = 10$ Ratern ist ein 95\%-KI von ungef\"ahr $\pm 0{,}15$ zu erwarten \citep{Shrout1979}. Eine Erh\"ohung auf $k = 30$ Vignetten w\"urde die Pr\"azision substanziell verbessern und sollte bei ausreichenden Ressourcen angestrebt werden.

\textbf{Falsifikationskriterium:} Falls $\kappa < 0{,}6$ oder ICC $< 0{,}65$, ist die Inter-Rater-Reliabilit\"at unzureichend --- dies erfordert Revision der Achsenstruktur oder klarere Operationalisierungen.

\item[Phase~3: Klinische Felderprobung (N=200+)] Prospektive Anwendung in klinischen Settings (ambulant und station\"ar). Messung von: Interrater-Reliabilit\"at, konvergente Validit\"at gegen\"uber etablierten Systemen, inkrementelle Validit\"at der Abdeckungsanalyse (werden durch Ii/IIIi identifizierte L\"ucken klinisch best\"atigt?), Akzeptanz und Usability (System Usability Scale, SUS).
\item[Phase~4: Multi-Site-Trial] Multizentrische Studie zur Generalisierbarkeit. Vergleich station\"ar vs.\ ambulant, Erwachsene vs.\ Kinder/Jugendliche, verschiedene kulturelle Kontexte.
\end{description}

Der MentalBench-Benchmark \citep{Song2026MentalBench} zeigt, dass gro\ss{}e Sprachmodelle systematisch an komplexen Dauer-, Funktions- und Ausschlusskriterien in der psychiatrischen Diagnostik scheitern, was die Begr\"undung f\"ur regelbasierte Entscheidungsunterst\"utzungssysteme mit expliziter Gatekeeper-Logik gegen\"uber rein statistischen Ans\"atzen verst\"arkt. Angesichts der rapiden Fortschritte in der LLM-basierten klinischen Unterst\"utzung ist zu betonen, dass das vorliegende System und KI-basierte Ans\"atze komplement\"ar, nicht konkurrierend sind: Das 6-Achsen-Modell liefert die strukturierte Architektur und die explizite Entscheidungslogik, die f\"ur klinische Nachvollziehbarkeit und Haftungssicherheit unerl\"asslich sind --- Eigenschaften, die Black-Box-Modelle systemisch nicht bieten k\"onnen. Zuk\"unftige Implementierungen k\"onnten LLMs als \textit{Assistenzsysteme innerhalb} der regelbasierten Architektur nutzen (z.\,B.\ f\"ur Freitext-Analyse von Anamnese-Berichten), ohne die Transparenz des Entscheidungspfads zu kompromittieren.


% ===================================================================
% 12. INTEROPERABILITÄT
% ===================================================================
\section{Interoperabilit\"at und Gesundheits-IT-Standards}
\label{sec:interop}

F\"ur die Integration in bestehende Gesundheitsinformationssysteme muss das 6-Achsen-Modell standardkonforme Schnittstellen implementieren.

\textbf{HL7 FHIR (R4):} Die Achsen lassen sich auf FHIR-Ressourcen abbilden: Achse~I/III-Diagnosen als \texttt{Condition}-Ressourcen mit Extensions f\"ur Subachsen-Zuordnung; Screening-Ergebnisse als \texttt{QuestionnaireResponse}; Funktionsstatus (Achse~IV) als \texttt{ClinicalImpression}; die Fallformulierung (Achse~V) als \texttt{CarePlan}; Belege (Achse~VI) als \texttt{DocumentReference}.

\textbf{SNOMED CT:} Alle Diagnosen sollten neben DSM-5-TR- und ICD-11-Codes auch SNOMED-CT-Konzept-IDs f\"uhren, um internationale Interoperabilit\"at zu gew\"ahrleisten.

\textbf{LOINC:} F\"ur Laborwerte (Achse~IIIe/IIIf), klinische Beobachtungen und Screening-Ergebnisse sollten LOINC-Codes (Logical Observation Identifiers Names and Codes) verwendet werden. Dies erm\"oglicht die standardisierte Dokumentation und den Austausch von Laborbefunden und psychometrischen Testergebnissen in der Evidenz-Matrix (Achse~VI).

\textbf{MHIRA-Kompatibilit\"at:} Das MHIRA-Projekt (Mental Health Information Reporting Assistant, BMC Psychiatry 2023) stellt die relevanteste Open-Source-Referenzarchitektur dar: ein cloud-basiertes, Docker-deployiertes psychiatrisches EHR-System mit digitalisierten psychometrischen Instrumenten. Eine FHIR-basierte Schnittstelle zum MHIRA-System sollte priorisiert werden.


% ===================================================================
% 13. DISKUSSION UND AUSBLICK
% ===================================================================
\section{Diskussion und Ausblick}
\label{sec:diskussion}

Das vorgestellte 6-Achsen-Modell ist kein R\"uckgriff auf das DSM-IV, sondern repr\"asentiert eine \textbf{vorgeschlagene Weiterentwicklung}, auf die die psychiatrische Fachwelt selbst konvergiert. Zehn vorgeschlagene Designmerkmale verdienen besondere Hervorhebung:

\textbf{Erstens} erweitert die Abdeckungsanalyse (Abschnitt~\ref{sec:abdeckung}) bestehende Entscheidungsunterst\"utzungsmethoden um eine quantitative Qualit\"atssicherungsmetrik. W\"ahrend verwandte Ans\"atze in Bayesianischen Diagnosesystemen existieren, liegt der spezifische Beitrag in der klinischen Qualit\"atssicherungsperspektive: Sie adressiert das von \citet{PlanaRipoll2019} dokumentierte Grundproblem transdiagnostischer Komorbidit\"at und wird symmetrisch in Achse~I (Ii) und Achse~III (IIIi) implementiert.

\textbf{Zweitens} verbr\"uckt die operationalisierte Integration von Fallformulierung (Achse~V) mit diagnostischer Klassifikation die langj\"ahrige Kluft zwischen \glqq Welche St\"orung hat dieser Patient?\grqq{} und \glqq Warum hat dieser Patient diese St\"orung?\grqq{}.

\textbf{Drittens} adressiert die Dual-System-Architektur DSM-5-TR/ICD-11 ein reales klinisches Bed\"urfnis in Kontexten, in denen beide Systeme Anwendung finden.

\textbf{Viertens} stellt die symmetrische Parallelstruktur von Achse~I und Achse~III sicher, dass das System als gemeinsames Werkzeug f\"ur alle beteiligten Professionen funktioniert. In den hier verglichenen computergest\"utzten Diagnosesystemen (OPCRIT, DTREE, CATEGO, CIDI, MHIRA) wurde ein solches Designprinzip nicht identifiziert; eine systematische Literaturrecherche, die auch weniger bekannte oder institutionsinterne Systeme umfasst, steht jedoch noch aus und ist f\"ur die endg\"ultige Bewertung der Neuheit dieses Ansatzes erforderlich.

\textbf{F\"unftens} bietet die Erweiterung auf 11~St\"orungsmodule eine vollst\"andige Abdeckung zwischen Screening und Diagnostik: Jedes Screening-Instrument m\"undet in einen strukturierten diagnostischen Workflow.

\textbf{Sechstens: Anwendung in Ausbildung und Supervision.} Das System eignet sich besonders als Lehr- und Supervisionswerkzeug in der psychiatrischen und psychotherapeutischen Ausbildung. Die PRO/CONTRA-Struktur zwingt Auszubildende, ihre diagnostischen Entscheidungen explizit zu begr\"unden und Gegenargumente systematisch zu dokumentieren --- eine F\"ahigkeit, die in konventioneller Freitext-Diagnostik nicht strukturell gefordert wird. Der Diagnostic Coverage Index (DCI) liefert Supervisoren eine quantitative R\"uckmeldung \"uber die diagnostische Vollst\"andigkeit der Auszubildenden: Ein niedriger DCI-Wert macht diagnostische L\"ucken sichtbar, die in Fallbesprechungen adressiert werden k\"onnen. In der psychotherapeutischen Ausbildung bietet Achse~V (Bedingungsmodell) eine strukturierte Alternative zu unstrukturierten Fallformulierungen, die h\"aufig idiosynkratisch und schwer vergleichbar sind.

\textbf{Siebtens} operationalisiert die PRO/CONTRA-Evidenzbewertung mit numerischer Konfidenz die differenzialdiagnostische Transparenz: Der Kliniker dokumentiert nicht nur, \textit{welche} Diagnose er stellt, sondern auch \textit{warum} --- einschlie\ss{}lich der bewusst verworfenen Gegenargumente. Diese Transparenz erm\"oglicht \textit{Shared Decision Making} (SDM): Patienten k\"onnen nachvollziehen, auf welchen Befunden ihre Diagnose basiert und welche Unsicherheiten bestehen, was die therapeutische Allianz und die Patientenautonomie f\"ordert. Die strukturierte PRO/CONTRA-Darstellung entspricht dabei dem dritten Schritt des SDM-Modells (Informationsaustausch und Deliberation) und schafft eine gemeinsame Entscheidungsgrundlage zwischen Kliniker und Patient.

\textbf{Achtens} implementiert das CAVE-Warnsystem in Achse~VI eine achsen\"ubergreifende Sicherheitsschicht, die kritische klinische Risiken (Labor-Artefakte, Medikamenten-Interaktionen, Kontraindikationen, zeitliche Fehlzuordnungen) prominent kennzeichnet. Angesichts der bekannten Polypharmazie-Problematik in der Psychiatrie --- insbesondere bei komorbiden somatischen Erkrankungen --- adressiert das CAVE-System eine h\"aufige klinische Fehlerquelle, die in bestehenden Klassifikationssystemen nicht systematisch ber\"ucksichtigt wird.

\textbf{Neuntens} erg\"anzt der longitudinale Symptomverlauf die querschnittliche Diagnostik um eine zeitliche Dimension: Die Dokumentation von Symptombeginn, aktuellem Status und Therapieansprechen erm\"oglicht die Identifikation von Chronifizierungsmustern und Therapieresistenz, die f\"ur die Behandlungsplanung essentiell sind. Achse~VI realisiert damit de facto ein \textit{integriertes Routine Outcome Monitoring (ROM)}-System im Sinne von \citet{Barkham2023}: Die regelm\"a\ss{}ige Erfassung validierter Instrumente (PHQ-9, GAD-7, PCL-5) mit Verlaufsvisualisierung erm\"oglicht die fr\"uhzeitige Erkennung von Not-on-Track-Verl\"aufen und unterst\"utzt adaptive Behandlungsplanung.

\textbf{Zehntens} realisiert die Designphilosophie des \textit{lebenden Dokuments} ein Paradigma, das die diagnostische Akte nicht als einmalig auszuf\"ullendes Formular, sondern als mitwachsende Prozessdokumentation versteht. Bezugspersonen-Netzwerk (Achse~IV), pr\"agende Erfahrungen und Grundkonflikte (Achse~II) sowie das Kontakt- und Beobachtungsprotokoll (Achse~VI) sind offene Listen, die den diagnostischen Prozess selbst dokumentieren und kontinuierlich Material f\"ur die Fallformulierung (Achse~V) liefern.

\vspace{0.5em}
\textbf{Methodische Einordnung:} Im Design-Science-Rahmen \citep{Hevner2004, Peffers2007} umfasst der vorliegende Artikel die Phasen Problem Identification, Design \& Development sowie Demonstration (Prototyp und Fallvignette). Die formale Evaluationsphase --- die f\"ur den vollst\"andigen Design-Science-Zyklus konstitutiv ist --- steht noch aus und wird durch die in Abschnitt~\ref{sec:validierung} beschriebene 4-Phasen-Strategie adressiert. Der Artikel versteht sich daher als Beitrag zur Designforschung, nicht als abgeschlossene Evaluation.

\textbf{Klinische Praktikabilit\"at und Bearbeitungszeit:} Eine zentrale Frage f\"ur die klinische Adoption ist die Bearbeitungszeit. Bei vollst\"andiger Ersterhebung aller sechs Achsen ist mit einer Bearbeitungszeit von 90--120~Minuten zu rechnen (einschlie\ss{}lich Screening, Anamnese und Dokumentation) --- vergleichbar mit einer umfassenden psychiatrischen Erstaufnahme, aber deutlich l\"anger als eine fokussierte ambulante Konsultation (20--50~Minuten). Das Skelett-Prinzip erm\"oglicht jedoch eine gestufte Nutzung: Eine Minimalerhebung (Achse~I mit Cross-Cutting-Screening und Achse~IV mit WHODAS~2.0) ist in 30--45~Minuten m\"oglich; die \"ubrigen Achsen werden im Verlauf erg\"anzt. Die Differenzierung ambulanter und station\"arer Nutzung ist f\"ur die geplante Validierung essentiell: Station\"are Teams k\"onnen die Achsen arbeitsteilig und \"uber mehrere Tage bearbeiten; ambulante Einzelbehandler ben\"otigen den gestuften Ansatz. Diese Unterscheidung sollte in Phase~3 der Validierung (klinische Felderprobung) systematisch untersucht werden.

\textbf{Verh\"altnis zur Process-Based Therapy (PBT):} Die PBT-Bewegung argumentiert gegen diagnostische Kategorien und f\"ur eine rein prozessbasierte Interventionsplanung. Das 6-Achsen-Modell positioniert sich als Br\"ucke: Es beh\"alt kategoriale Diagnostik bei (klinisch, administrativ und versicherungsrechtlich notwendig), integriert aber \"uber Achse~V (Bedingungsmodell) und die Abdeckungsanalyse prozessorientierte Elemente --- insbesondere die Identifikation aufrechterhaltender Faktoren (P3) und ungekl\"arter Symptome, die auf behandelbare Prozesse jenseits der Kategorienlogik hinweisen k\"onnen.

Die technische Implementierung sollte vier Meilensteine priorisieren: (1)~das Cross-Cutting-Screening-Modul als Einstiegspunkt des Systems, (2)~die Gatekeeper-Zustandsmaschine mittels \texttt{transitions} HSM, (3)~das PID-5-BF-Assessment mit Plotly-Radardiagramm und WeasyPrint-PDF-Export, und (4)~die FHIR-Schnittstellendefinition f\"ur Interoperabilit\"at.

\textbf{Implementierungsbarrieren:} Die Erfahrung mit fr\"uheren multiaxialen Systemen und klinischen IT-Einf\"uhrungen zeigt, dass technische Funktionalit\"at allein keine Adoption gew\"ahrleistet. Zu den erwartbaren Barrieren geh\"oren: (1)~Zeitmangel im klinischen Alltag, insbesondere in unterbesetzten ambulanten Settings; (2)~institutionelle Tr\"agheit und Widerstand gegen neue Dokumentationssysteme, die zun\"achst als Mehrbelastung wahrgenommen werden; (3)~die Notwendigkeit interdisziplin\"arer Abstimmung, die in fragmentierten Versorgungsstrukturen organisatorisch anspruchsvoll ist; und (4)~die Integration in bestehende Krankenhausinformationssysteme (KIS), die propriet\"are Schnittstellen und eingeschr\"ankte Erweiterbarkeit aufweisen. Das Skelett-Prinzip und die gestufte Nutzung (vgl.\ Bearbeitungszeit oben) adressieren die Zeitbarriere partiell; die \"ubrigen Barrieren erfordern eine begleitende Implementierungsstrategie im Sinne des CFIR-Frameworks \citep{Damschroder2009}.


% ===================================================================
% 14. LIMITATIONEN UND ZUKÜNFTIGER FORSCHUNGSBEDARF
% ===================================================================
\section{Limitationen und zuk\"unftiger Forschungsbedarf}
\label{sec:limitationen}

Das vorgestellte 6-Achsen-Modell unterliegt mehreren Limitationen, die bei der Interpretation und Weiterentwicklung ber\"ucksichtigt werden m\"ussen.

\textbf{Fehlende empirische Validierung:} Die zentrale Limitation besteht darin, dass das Modell bislang rein konzeptionell-theoretisch entwickelt wurde. Weder Interrater-Reliabilit\"at noch konvergente oder inkrementelle Validit\"at sind empirisch gepr\"uft. Die in Abschnitt~\ref{sec:validierung} beschriebene 4-Phasen-Validierungsstrategie stellt einen Forschungsplan dar, keine vorliegenden Ergebnisse. Bis zur Durchf\"uhrung mindestens der Phasen~1 und~2 bleibt der klinische Nutzen des Modells eine begr\"undete Hypothese.

\textbf{Komplexit\"at und klinische Praktikabilit\"at:} Das System umfasst sechs Achsen mit insgesamt 40~Subachsen. Obwohl das Skelett-Prinzip keine Pflichtfelder vorsieht, k\"onnte die Komplexit\"at in der klinischen Praxis abschreckend wirken. Eine wesentliche Forschungsfrage ist, ob Kliniker das System tats\"achlich in der intendierten Tiefe nutzen oder auf wenige Achsen reduzieren. Usability-Studien (SUS-Skala) sind essenziell.

\textbf{Skalierbarkeit der Gewichtsmatrix:} Die Abdeckungsanalyse setzt eine vollst\"andige Gewichtsmatrix $w(s_i, d_j)$ f\"ur alle Symptom-Diagnose-Kombinationen voraus. Bei \"uber 300 DSM-5-TR-Diagnosen und mehreren hundert ICD-11-Codes ist die manuelle Erstellung und Pflege dieser Matrix ein erheblicher Aufwand. Eine pragmatische L\"osung besteht darin, die Gewichte nur f\"ur die im konkreten Fall relevanten Diagnosen zu berechnen (typischerweise 3--8 Differenzialdiagnosen), wodurch die Matrix auf handhabbare Gr\"o\ss{}e reduziert wird. Langfristig k\"onnten die Gewichte aus strukturierten Kriterienkatalogen (z.\,B.\ SNOMED-CT-Konzepte, ICD-11 Foundation Layer) semi-automatisch abgeleitet werden.

\textbf{Keine Normierung der Abdeckungsanalyse:} Die quantitative Abdeckungsanalyse (Achse~Ii/IIIi) verwendet prozentuale Metriken, f\"ur die keine empirischen Normen existieren. Die Schwellenwerte (85\%/60\%) sind durch klinische Analogien motiviert (vgl.\ Abschnitt~\ref{sec:abdeckung}), aber nicht empirisch validiert. Die Sensitivit\"atsanalyse zeigt Robustheit bei $n \geq 8$ Symptomen, aber Instabilit\"at bei kleinen Symptomzahlen. Die empirische Kalibrierung in mindestens drei klinischen Kohortengruppen (station\"ar, ambulant, somatisch komorbid) ist eine Voraussetzung f\"ur klinische Adoption.

\textbf{Eingeschr\"ankte Literaturbasis:} Als Designarbeit st\"utzt sich diese Arbeit prim\"ar auf diagnostische Klassifikationshandb\"ucher und eine kleine Anzahl grundlegender Publikationen. Ein systematisches Review aller existierenden computergest\"utzten Diagnosesysteme (z.\,B.\ OPCRIT, DTREE, CATEGO) und deren empirische Erfolgsbilanz w\"urde die Begr\"undung f\"ur die hier getroffenen spezifischen Designentscheidungen st\"arken.

\textbf{Dokumentation des KI-gest\"utzten Prozesses:} Der substanzielle Einsatz KI-gest\"utzter Werkzeuge sowohl am Systemdesign als auch an der Manuskripterstellung (d.\,h.\ Literaturorganisation, Textstrukturierung, \"Ubersetzungsunterst\"utzung, Konsistenzpr\"ufung und Formulierungshilfe; vgl.\ Methoden, Phase~4) wirft Fragen zur Prozesstransparenz auf. W\"ahrend der Quellcode \"offentlich verf\"ugbar ist, ist die iterative Prompt-, Review- und Revisionsgeschichte --- einschlie\ss{}lich der Expertenfeedback-Zyklen --- nicht vollst\"andig so dokumentiert, dass eine vollst\"andige Prozessrekonstruktion m\"oglich w\"are.

\textbf{Kulturelle Generalisierbarkeit und jurisdiktionsspezifische Module:} Das Modell integriert das Cultural Formulation Interview (CFI), wurde aber prim\"ar aus westlich-europ\"aischer Perspektive entwickelt. Jurisdiktionsspezifische Subachsen wie der GdB (Achse~IVb) sind als austauschbare Module konzipiert: F\"ur internationale Implementierungen sollte der GdB durch das jeweilige nationale Behinderungsbewertungssystem ersetzt werden (z.\,B.\ Disability Adjusted Life Years in WHO-Kontexten, Social Security Disability Determination in den USA). Die Anwendbarkeit in nicht-westlichen klinischen Kontexten bedarf gesonderter Pr\"ufung. Insbesondere die Pers\"onlichkeitsdiagnostik mittels PID-5 und die biographischen Komponenten (Achse~II) k\"onnen kulturell unterschiedlich interpretiert werden.

\textbf{Technologieabh\"angigkeit:} Die Implementierung als computergest\"utztes Entscheidungsunterst\"utzungssystem setzt eine funktionsf\"ahige IT-Infrastruktur voraus. In ressourcenarmen Settings oder bei Systemausf\"allen muss das Modell auch als papierbasiertes Schema funktionieren --- eine Anforderung, die bislang nicht systematisch gepr\"uft wurde.

\textbf{Fehlende Stakeholder-Perspektive und Kosteneffektivit\"at:} Der vorliegende Artikel ist prim\"ar aus der Systementwickler-Perspektive geschrieben. Wesentliche Perspektiven fehlen: Wie erleben Patienten eine derart umfassende diagnostische Erhebung? Wie verh\"alt sich die Bearbeitungszeit im Vergleich zur konventionellen Diagnostik? Wie bewerten Kostentr\"ager (z.\,B.\ Krankenkassen) den zus\"atzlichen diagnostischen Aufwand? Zuk\"unftige Studien sollten eine Bearbeitungszeitanalyse, eine Kosten-Nutzen-Bewertung und eine systematische Erhebung der Patientenperspektive einschlie\ss{}en.

\textbf{Learning-Health-System-Perspektive:} Die strukturierten, standardisiert kodierten Daten des 6-Achsen-Modells (ICD-11, SNOMED~CT, ICF, PID-5) eignen sich ideal f\"ur sekund\"are Forschungsnutzung. In einem Learning-Health-System-Paradigma tr\"agt jede diagnostische Episode zur iterativen Verbesserung der Gewichtsmatrix der Abdeckungsanalyse bei: Empirische Symptom-Diagnose-Assoziationen aus aggregierten, anonymisierten Patientendaten k\"onnen die initialen Gewichte schrittweise durch datengest\"utzte Sch\"atzungen ersetzen. Diese Perspektive verbindet die klinische Routineanwendung mit der kontinuierlichen Systemoptimierung.

\textbf{Zuk\"unftiger Forschungsbedarf:} Neben der empirischen Validierung (Abschnitt~\ref{sec:validierung}) sind vier Forschungsrichtungen priorit\"ar: (1)~die Entwicklung eines abgek\"urzten \glqq Rapid Assessment\grqq{}-Modus f\"ur Notaufnahmen und Erstgespr\"ache, (2)~die Integration pr\"adiktiver Modelle (maschinelles Lernen) zur Unterst\"utzung der Differenzialdiagnostik auf Basis der akkumulierten Patientendaten, (3)~l\"angsschnittliche Studien zur Sensitivit\"at des Symptomverlaufs (Achse~VI) f\"ur Therapieansprechen und R\"uckfallpr\"adiktion, und (4)~die Anbindung an Ecological Momentary Assessment (EMA) und Digital Phenotyping: Smartphone-basierte Echtzeit-Datenerhebung k\"onnte den longitudinalen Symptomverlauf (Achse~VI) mit kontinuierlichen Verhaltensdaten anreichern und die zeitliche Aufl\"osung der diagnostischen Dokumentation substanziell erh\"ohen. Dar\"uber hinaus sollte die tats\"achliche klinische Einf\"uhrung durch einen Implementation-Science-Rahmen --- etwa das Consolidated Framework for Implementation Research (CFIR) --- begleitet werden, der Barrieren und F\"orderfaktoren der Adoption in den f\"unf CFIR-Dom\"anen (Intervention Characteristics, Outer Setting, Inner Setting, Individuals, Process) systematisch erfasst. Die historische Erfahrung mit dem DSM-IV-Multiaxialsystem zeigt, dass selbst gut konzipierte Systeme an inkonsistenter klinischer Nutzung scheitern k\"onnen.


% ===================================================================
% 15. ZUSAMMENFASSUNG
% ===================================================================
\section{Zusammenfassung}
\label{sec:zusammenfassung}

Das vorgeschlagene 6-Achsen-System adressiert eine strukturelle L\"ucke in der psychiatrischen Diagnostik, die seit der Abschaffung des DSM-IV-Multiaxialsystems besteht. F\"ur die klinische Praxis liegt der prim\"are Beitrag des Modells in der Wiederherstellung --- und substanziellen Erweiterung --- der strukturierten Aufforderung zur umfassenden biopsychosozialen Beurteilung. Die symmetrische Achse-I/III-Architektur erm\"oglicht strukturierte interprofessionelle Zusammenarbeit, die Abdeckungsanalyse liefert Klinikern eine transparente, quantitative Metrik f\"ur diagnostische Vollst\"andigkeit, und das CAVE-Warnsystem adressiert die wachsende Herausforderung des achsen\"ubergreifenden Risikomanagements bei komplexen F\"allen. Die PRO/CONTRA-Evidenzbewertung mit Konfidenzsch\"atzung operationalisiert diagnostische Transparenz in einer Weise, die sowohl klinische Rechenschaftspflicht als auch Ausbildung unterst\"utzt.

Die Architektur des Modells ist konsistent mit der institutionellen Entwicklungsrichtung (vgl.\ Abschnitt~\ref{sec:einleitung}). Die Konvergenz mit der expliziten Forderung von \citet{ErlichFirst2025} nach einer R\"uckkehr zur strukturierten axialen Beurteilung unterstreicht die Aktualit\"at dieses Beitrags. Zugleich geht das Modell \"uber institutionelle Revisionen hinaus, insbesondere durch die Abdeckungsanalyse und die operationalisierte Fallformulierung (Achse~V), die Klassifikation und Behandlungsplanung verbr\"uckt.

Konkrete n\"achste Schritte umfassen: (1)~eine fokussierte Expertenbegutachtung ($N = 5$--$10$) der Achsenstruktur als unmittelbarer Minimalschritt, (2)~die Interrater-Reliabilit\"atsstudie mit standardisierten Fallvignetten (Phase~2, Abschnitt~\ref{sec:validierung}), (3)~die Entwicklung der HL7-FHIR-Mapping-Spezifikation, (4)~die Publikation der formalen mathematischen Spezifikation der Abdeckungsanalyse und (5)~empirische Vergleiche von Freitext-Dokumentation versus strukturierter multiaxialer Diagnostik. Der vollst\"andige Quellcode ist \"offentlich verf\"ugbar, um unabh\"angige Evaluation und Replikation zu erm\"oglichen.

Zusammengefasst: Das 6-Achsen-Modell versteht sich als Entwurf einer diagnostischen Architektur, die das Beste aus kategorialer Klassifikation, dimensionaler Erfassung und individualisierter Fallformulierung in einem einzigen, computergest\"utzten System vereint --- ein Entwurf, der seine klinische Bew\"ahrungsprobe noch vor sich hat.


% ===================================================================
% LITERATURVERZEICHNIS
% ===================================================================
\begin{thebibliography}{99}
\bibitem[Aboraya et al.(2006)]{Aboraya2006}
Aboraya, A., Rankin, E., France, C., El-Missiry, A. \& John, C. (2006).
\newblock The reliability of psychiatric diagnosis revisited: The clinician's guide to improve the reliability of psychiatric diagnosis.
\newblock \textit{Psychiatry (Edgmont)}, 3(1), 41--50. PMID: 21103149.

\bibitem[Achenbach \& Rescorla(2001)]{Achenbach2001}
Achenbach, T. M. \& Rescorla, L. A. (2001).
\newblock \textit{Manual for the ASEBA School-Age Forms \& Profiles}.
\newblock Burlington, VT: University of Vermont, Research Center for Children, Youth, \& Families.

\bibitem[Allison et al.(2012)]{Allison2012}
Allison, C., Auyeung, B. \& Baron-Cohen, S. (2012).
\newblock Toward brief ``red flags'' for autism screening: The Short Autism Spectrum Quotient and the Short Quantitative Checklist in 1,000 cases and 3,000 controls.
\newblock \textit{Journal of the American Academy of Child \& Adolescent Psychiatry}, 51(2), 202--212.e7. DOI: 10.1016/j.jaac.2011.11.003.

\bibitem[APA(1980)]{APA1980DSM3}
American Psychiatric Association (1980).
\newblock \textit{Diagnostic and Statistical Manual of Mental Disorders} (3rd ed.).
\newblock Washington, DC: American Psychiatric Publishing.

\bibitem[APA(2013)]{APA2013}
American Psychiatric Association (2013).
\newblock \textit{Diagnostic and Statistical Manual of Mental Disorders} (5th ed.).
\newblock Arlington, VA: American Psychiatric Publishing.

\bibitem[APA(2022)]{APA2022}
American Psychiatric Association (2022).
\newblock \textit{Diagnostic and Statistical Manual of Mental Disorders, Text Revision} (DSM-5-TR).
\newblock Washington, DC: American Psychiatric Publishing.

\bibitem[Ayuso-Mateos et al.(2013)]{AyusoMateos2013Bipolar}
Ayuso-Mateos, J.\,L., Avila, C.\,C., Anaya, C., Cieza, A. \& Vieta, E. (2013).
\newblock Development of the International Classification of Functioning, Disability and Health core sets for bipolar disorders: Results of an international consensus process.
\newblock \textit{Disability and Rehabilitation}, 35(25), 2138--2146. DOI: 10.3109/09638288.2013.771708.

\bibitem[Bach et al.(2022)]{Bach2022}
Bach, B., Kramer, U., Doering, S. et al. (2022).
\newblock The ICD-11 classification of personality disorders: A European perspective on challenges and opportunities.
\newblock \textit{Borderline Personality Disorder and Emotion Dysregulation}, 9(1), 12.

\bibitem[Barkham et al.(2023)]{Barkham2023}
Barkham, M., De Jong, K., Delgadillo, J. \& Lutz, W. (2023).
\newblock Routine Outcome Monitoring (ROM) and Feedback: Research Review and Recommendations.
\newblock \textit{Psychotherapy Research}, 33(7), 841--855. DOI: 10.1080/10503307.2023.2181114.

\bibitem[Barnhill(2014)]{BarnhillDSM5ClinicalCases}
Barnhill, J.\,W. (Ed.) (2014).
\newblock \textit{DSM-5 Clinical Cases}.
\newblock Washington, DC: American Psychiatric Publishing.

\bibitem[Mayall et al.(2024)]{BJPsychAdv2024}
Mayall, M., McDermott, B., Sadhu, R., Teoh, Y., Bosanquet, M. \& Nundeekasen, S. (2024).
\newblock Practical aspects of multiaxial classification: A clinically useful biopsychosocial framework for child and adolescent psychiatry.
\newblock \textit{BJPsych Advances}, 30(4), 242--256. DOI: 10.1192/bja.2023.39.

\bibitem[Blevins et al.(2015)]{Blevins2015}
Blevins, C.\,A., Weathers, F.\,W., Davis, M.\,T. et al. (2015).
\newblock The Posttraumatic Stress Disorder Checklist for DSM-5 (PCL-5).
\newblock \textit{Journal of Traumatic Stress}, 28(6), 489--498.

\bibitem[Bolton(2008)]{Bolton2008}
Bolton, D. (2008).
\newblock \textit{What is Mental Disorder? An Essay in Philosophy, Science, and Values}.
\newblock Oxford: Oxford University Press.

\bibitem[Carlson \& Putnam(1993)]{CarlsonPutnam1993}
Carlson, E.\,B. \& Putnam, F.\,W. (1993).
\newblock An update on the Dissociative Experiences Scale.
\newblock \textit{Dissociation}, 6(1), 16--27.

\bibitem[Cieza et al.(2004)]{Cieza2004Depression}
Cieza, A., Chatterji, S., Andersen, C. et al. (2004).
\newblock ICF Core Sets for depression.
\newblock \textit{Journal of Rehabilitation Medicine}, 36(Suppl. 44), 128--134. DOI: 10.1080/16501960410016055.

\bibitem[Cloitre et al.(2018)]{Cloitre2018}
Cloitre, M., Shevlin, M., Brewin, C.\,R. et al. (2018).
\newblock The International Trauma Questionnaire: Development of a self-report measure of ICD-11 PTSD and Complex PTSD.
\newblock \textit{Acta Psychiatrica Scandinavica}, 138(6), 536--546.

\bibitem[Damschroder et al.(2009)]{Damschroder2009}
Damschroder, L.\,J., Aron, D.\,C., Keith, R.\,E. et al. (2009).
\newblock Fostering implementation of health services research findings into practice: A consolidated framework for advancing implementation science.
\newblock \textit{Implementation Science}, 4, 50.

\bibitem[Erlich et al.(2025)]{ErlichFirst2025}
Erlich, M.\,D., First, M.\,B., Talley, R.\,M. et al. (2025).
\newblock Integrating social determinants of health into clinical formulations: The need for a biaxial system in the DSM and psychiatry.
\newblock \textit{Psychiatric Services}, 76(10), 911--914.

\bibitem[First et al.(1993)]{First1993}
First, M.\,B., Williams, J.\,B.\,W. \& Spitzer, R.\,L. (1993).
\newblock DTREE: A decision-support system for DSM-III-R and DSM-IV diagnoses.
\newblock \textit{Psychopharmacology Bulletin}, 29(3), 255--263.

\bibitem[First(2024)]{First2024}
First, M.\,B. (2024).
\newblock \textit{DSM-5-TR Handbook of Differential Diagnosis}.
\newblock Washington, DC: American Psychiatric Publishing.

\bibitem[Foa et al.(2002)]{Foa2002}
Foa, E.\,B., Huppert, J.\,D., Leiberg, S. et al. (2002).
\newblock The Obsessive-Compulsive Inventory: Development and validation of a short version.
\newblock \textit{Psychological Assessment}, 14(4), 485--496.

\bibitem[Gierk et al.(2014)]{Gierk2014}
Gierk, B., Kohlmann, S., Kroenke, K., Spangenberg, L., Zenger, M., Br\"ahler, E. \& L\"owe, B. (2014).
\newblock The Somatic Symptom Scale-8 (SSS-8): A brief measure of somatic symptom burden.
\newblock \textit{JAMA Internal Medicine}, 174(3), 399--407.

\bibitem[G{\'o}mez-Benito et al.(2018)]{GomezBenito2018Schizophrenia}
G{\'o}mez-Benito, J., Guilera, G., Barrios, M. et al. (2018).
\newblock Beyond diagnosis: The Core Sets for persons with schizophrenia based on the World Health Organization's International Classification of Functioning, Disability, and Health.
\newblock \textit{Disability and Rehabilitation}, 40(23), 2756--2766. DOI: 10.1080/09638288.2017.1356384.

\bibitem[Gspandl et al.(2018)]{Gspandl2018}
Gspandl, S., Peirson, R.\,P., Nahhas, R.\,W., Skale, T.\,G. \& Lehrer, D.\,S. (2018).
\newblock Comparing Global Assessment of Functioning (GAF) and World Health Organization Disability Assessment Schedule (WHODAS) 2.0 in schizophrenia.
\newblock \textit{Psychiatry Research}, 259, 251--253.

\bibitem[Hevner et al.(2004)]{Hevner2004}
Hevner, A.\,R., March, S.\,T., Park, J. \& Ram, S. (2004).
\newblock Design science in information systems research.
\newblock \textit{MIS Quarterly}, 28(1), 75--105.

\bibitem[Insel et al.(2010)]{Insel2010}
Insel, T., Cuthbert, B., Garvey, M. et al. (2010).
\newblock Research Domain Criteria (RDoC): Toward a new classification framework for research on mental disorders.
\newblock \textit{American Journal of Psychiatry}, 167(7), 748--751.

\bibitem[Ising et al.(2012)]{Ising2012}
Ising, H.\,K., Veling, W., Loewy, R.\,L. et al. (2012).
\newblock The validity of the 16-item version of the Prodromal Questionnaire (PQ-16) to screen for ultra high risk of developing psychosis.
\newblock \textit{Schizophrenia Research}, 138(2--3), 207--212.

\bibitem[Keeley et al.(2016)]{Keeley2016}
Keeley, J.\,W., Reed, G.\,M., Roberts, M.\,C. et al. (2016).
\newblock Developing a science of clinical utility in diagnostic classification systems: Field study strategies for ICD-11 mental and behavioural disorders.
\newblock \textit{American Psychologist}, 71(1), 3--16.

\bibitem[Kessler and \"{U}st\"{u}n(2004)]{Kessler2004}
Kessler, R.\,C. \& \"{U}st\"{u}n, T.\,B. (2004).
\newblock The World Mental Health (WMH) Survey Initiative version of the World Health Organization (WHO) Composite International Diagnostic Interview (CIDI).
\newblock \textit{International Journal of Methods in Psychiatric Research}, 13(2), 93--121.

\bibitem[Kessler et al.(2005)]{Kessler2005}
Kessler, R.\,C., Adler, L., Ames, M. et al. (2005).
\newblock The World Health Organization Adult ADHD Self-Report Scale (ASRS).
\newblock \textit{Psychological Medicine}, 35(2), 245--256.

\bibitem[Kotov et al.(2017)]{Kotov2017}
Kotov, R., Krueger, R.\,F., Watson, D. et al. (2017).
\newblock The Hierarchical Taxonomy of Psychopathology (HiTOP): A dimensional alternative to traditional nosologies.
\newblock \textit{Journal of Abnormal Psychology}, 126(4), 454--477.

\bibitem[Kotov et al.(2021)]{Kotov2021}
Kotov, R., Krueger, R.\,F., Watson, D. et al. (2021).
\newblock The Hierarchical Taxonomy of Psychopathology (HiTOP): A quantitative nosology based on consensus of evidence.
\newblock \textit{Annual Review of Clinical Psychology}, 17, 83--108.

\bibitem[Kotov et al.(2022)]{Kotov2022}
Kotov, R., Cicero, D.\,C., Conway, C.\,C. et al. (2022).
\newblock The Hierarchical Taxonomy of Psychopathology (HiTOP) in psychiatric practice and research.
\newblock \textit{Psychological Medicine}, 52(9), 1666--1678. DOI: 10.1017/S0033291722001301.

\bibitem[Kress et al.(2014)]{Kress2014}
Kress, V.\,E., Barrio Minton, C.\,A., Adamson, N.\,A., Paylo, M.\,J. \& Pope, V. (2014).
\newblock The removal of the multiaxial system in the DSM-5: Implications and practice suggestions for counselors.
\newblock \textit{The Professional Counselor}, 4(3), 191--201.

\bibitem[Krueger et al.(2018)]{Krueger2018}
Krueger, R.\,F., Kotov, R., Watson, D. et al. (2018).
\newblock Progress in achieving quantitative classification of psychopathology.
\newblock \textit{World Psychiatry}, 17(3), 282--293.

\bibitem[Levis et al.(2019)]{Levis2019}
Levis, B., Benedetti, A. \& Thombs, B.\,D. (2019).
\newblock Accuracy of Patient Health Questionnaire-9 (PHQ-9) for screening to detect major depression.
\newblock \textit{BMJ}, 365, l1476.

\bibitem[Linden \& Baron(2005)]{LindenBaron2005}
Linden, M. \& Baron, S. (2005).
\newblock Das Mini-ICF-Rating f\"ur psychische St\"orungen (Mini-ICF-P): Ein Kurzinstrument zur Beurteilung von F\"ahigkeitsst\"orungen bei psychischen Erkrankungen.
\newblock \textit{Die Rehabilitation}, 44(3), 144--151. DOI: 10.1055/s-2004-834786.

\bibitem[Luciano et al.(2010)]{Luciano2010}
Luciano, J.\,V., Ayuso-Mateos, J.\,L., Aguado, J. et al. (2010).
\newblock The 12-item World Health Organization Disability Assessment Schedule II (WHO-DAS II): A nonparametric item response analysis.
\newblock \textit{BMC Medical Research Methodology}, 10, 45.

\bibitem[McGuffin et al.(1991)]{McGuffin1991}
McGuffin, P., Farmer, A. \& Harvey, I. (1991).
\newblock A polydiagnostic application of operational criteria in studies of psychotic illness: Development and reliability of the OPCRIT system.
\newblock \textit{Archives of General Psychiatry}, 48(8), 764--770.

\bibitem[Morgan et al.(1999)]{Morgan1999}
Morgan, J.\,F., Reid, F. \& Lacey, J.\,H. (1999).
\newblock The SCOFF questionnaire: Assessment of a new screening tool for eating disorders.
\newblock \textit{BMJ}, 319(7223), 1467--1468.

\bibitem[Morin et al.(2011)]{Morin2011}
Morin, C.\,M., Belleville, G., B\'elanger, L. \& Ivers, H. (2011).
\newblock The Insomnia Severity Index: Psychometric indicators to detect insomnia cases and evaluate treatment response.
\newblock \textit{Sleep}, 34(5), 601--608.

\bibitem[Narrow et al.(2013)]{Narrow2013}
Narrow, W.\,E., Clarke, D.\,E., Kuramoto, S.\,J. et al. (2013).
\newblock DSM-5 field trials in the United States and Canada, Part III.
\newblock \textit{American Journal of Psychiatry}, 170(1), 71--82.

\bibitem[Oltmanns \& Widiger(2018)]{OltmannsWidiger2018}
Oltmanns, J.\,R. \& Widiger, T.\,A. (2018).
\newblock A self-report measure for the ICD-11 dimensional trait model proposal: The Personality Inventory for ICD-11.
\newblock \textit{Psychological Assessment}, 30(2), 154--169.

\bibitem[Owen(2023)]{Owen2023}
Owen, G. (2023).
\newblock What is formulation in psychiatry?
\newblock \textit{Psychological Medicine}, 53(5), 1700--1707. DOI: 10.1017/S0033291723000016.

\bibitem[Peffers et al.(2007)]{Peffers2007}
Peffers, K., Tuunanen, T., Rothenberger, M.\,A. \& Chatterjee, S. (2007).
\newblock A design science research methodology for information systems research.
\newblock \textit{Journal of Management Information Systems}, 24(3), 45--77.

\bibitem[Plana-Ripoll et al.(2019)]{PlanaRipoll2019}
Plana-Ripoll, O., Pedersen, C.\,B., Holtz, Y. et al. (2019).
\newblock Exploring Comorbidity Within Mental Disorders Among a Danish National Population.
\newblock \textit{JAMA Psychiatry}, 76(3), 259--270.

\bibitem[Posner et al.(2011)]{Posner2011}
Posner, K., Brown, G.\,K., Stanley, B. et al. (2011).
\newblock The Columbia-Suicide Severity Rating Scale (C-SSRS).
\newblock \textit{American Journal of Psychiatry}, 168(12), 1266--1277.

\bibitem[Probst(2014)]{Probst2014}
Probst, B. (2014).
\newblock The life and death of Axis IV: Caught in the quest for a theory of mental disorder.
\newblock \textit{Research on Social Work Practice}, 24(1), 123--131. DOI: 10.1177/1049731513491326.

\bibitem[Regier et al.(2013)]{Regier2013}
Regier, D.\,A., Narrow, W.\,E., Clarke, D.\,E. et al. (2013).
\newblock DSM-5 field trials in the United States and Canada, Part II: Test-retest reliability of selected categorical diagnoses.
\newblock \textit{American Journal of Psychiatry}, 170(1), 59--70.

\bibitem[Riegel et al.(2021)]{Riegel2021}
Riegel, K.\,D., Ksinan, A.\,J. \& Schlosserova, L. (2021).
\newblock Psychometric properties of the independent 36-item PID5BF+M for ICD-11 in the Czech-speaking community sample.
\newblock \textit{Frontiers in Psychiatry}, 12, 643270. DOI: 10.3389/fpsyt.2021.643270.

\bibitem[Rosenman et al.(2003)]{Rosenman2003}
Rosenman, S.\,J., Korten, A.\,E., Medway, J. \& Evans, M. (2003).
\newblock Dimensional vs.\ categorical diagnosis in psychosis.
\newblock \textit{Acta Psychiatrica Scandinavica}, 107(5), 378--384.

\bibitem[Saunders et al.(1993)]{Saunders1993}
Saunders, J.\,B., Aasland, O.\,G., Babor, T.\,F. et al. (1993).
\newblock Development of the Alcohol Use Disorders Identification Test (AUDIT).
\newblock \textit{Addiction}, 88(6), 791--804.

\bibitem[Shrout \& Fleiss(1979)]{Shrout1979}
Shrout, P.\,E. \& Fleiss, J.\,L. (1979).
\newblock Intraclass correlations: Uses in assessing rater reliability.
\newblock \textit{Psychological Bulletin}, 86(2), 420--428.

\bibitem[Song et al.(2026)]{Song2026MentalBench}
Song, H., Kang, M., Shin, J. et al. (2026).
\newblock MentalBench: A DSM-Grounded Benchmark for Evaluating Psychiatric Diagnostic Capability of Large Language Models.
\newblock \textit{arXiv preprint arXiv:2602.12871}. DOI: 10.48550/arXiv.2602.12871.

\bibitem[Spitzer et al.(2006)]{Spitzer2006}
Spitzer, R.\,L., Kroenke, K., Williams, J.\,B.\,W. \& L\"owe, B. (2006).
\newblock A brief measure for assessing Generalized Anxiety Disorder: The GAD-7.
\newblock \textit{Archives of Internal Medicine}, 166(10), 1092--1097.

\bibitem[Zimmermann et al.(2023)]{Stasiak2023}
Zimmermann, R., Konjufca, J., Sakejo, P. et al. (2023).
\newblock Mental Health Information Reporting Assistant (MHIRA)---an open-source software facilitating evidence-based assessment for clinical services.
\newblock \textit{BMC Psychiatry}, 23, 706. DOI: 10.1186/s12888-023-05201-0.

\bibitem[Wakefield(1992)]{Wakefield1992}
Wakefield, J.\,C. (1992).
\newblock The concept of mental disorder: On the boundary between biological facts and social values.
\newblock \textit{American Psychologist}, 47(3), 373--388.

\bibitem[WHO(2010)]{WHO2010}
World Health Organization (2010).
\newblock \textit{WHODAS 2.0: WHO Disability Assessment Schedule 2.0}.
\newblock Geneva: WHO.

\bibitem[WHO(2019)]{WHO2019}
World Health Organization (2019).
\newblock \textit{International Classification of Diseases, 11th Revision} (ICD-11).
\newblock Geneva: WHO.

\bibitem[Williams(1985)]{Williams1985}
Williams, J.\,B.\,W. (1985).
\newblock The multiaxial system of DSM-III: Where did it come from and where should it go? I.\ Its origins and critiques.
\newblock \textit{Archives of General Psychiatry}, 42(2), 175--180.

\bibitem[Wing et al.(1974)]{Wing1974}
Wing, J.\,K., Cooper, J.\,E. \& Sartorius, N. (1974).
\newblock \textit{Measurement and Classification of Psychiatric Symptoms: An Instruction Manual for the PSE and CATEGO Program}.
\newblock Cambridge: Cambridge University Press.
\end{thebibliography}

\end{document}
